\documentclass[12pt,a4paper,openany]{report}
% ============================================================
% preambule.tex -- commun a tout le memoire.
% Compile avec pdflatex (Overleaf : babel francais) ou xelatex/lualatex (polyglossia).
% ============================================================
\usepackage{iftex}
\ifPDFTeX
  \usepackage[utf8]{inputenc}
  \usepackage[T1]{fontenc}
  \usepackage{mathpazo}
  \usepackage[french]{babel}
\else
  \usepackage{fontspec}
  \usepackage{polyglossia}
  \setmainlanguage{french}
  \setmainfont{TeX Gyre Pagella}
  \setmonofont{DejaVu Sans Mono}[Scale=0.85]
\fi
\usepackage{amsmath,amssymb,amsthm}
\usepackage{microtype}
\usepackage{graphicx}
\usepackage{booktabs,longtable,array,tabularx}
\usepackage{caption}
\usepackage{enumitem}
\usepackage{placeins}
\usepackage[table]{xcolor}
\usepackage{tikz}
\usetikzlibrary{arrows.meta,positioning,calc,patterns,decorations.pathreplacing}
\usepackage{pgfplots}\pgfplotsset{compat=1.16}
\usepackage[a4paper,margin=2.6cm]{geometry}
\usepackage{fancyhdr}
\usepackage{titlesec}
\emergencystretch=3em
\setlength{\parskip}{0.4em}
\setlength{\parindent}{0pt}

% ---- palette
\definecolor{navy}{HTML}{0B1D3A}
\definecolor{steel}{HTML}{3A6EA5}
\definecolor{ice}{HTML}{E7EEF6}
\definecolor{alert}{HTML}{D62839}
\definecolor{ok}{HTML}{2E8B57}
\definecolor{muted}{HTML}{6B7B8C}
\definecolor{amber}{HTML}{B7791F}

\usepackage[colorlinks=true,linkcolor=steel,citecolor=steel,urlcolor=steel]{hyperref}

\setcounter{secnumdepth}{3}\setcounter{tocdepth}{3}
\titleformat{\chapter}[display]{\normalfont\huge\bfseries\color{navy}}{\chaptertitlename\ \thechapter}{14pt}{\Huge}
\titleformat{name=\chapter,numberless}[block]{\normalfont\Large\scshape\color{navy}\raggedleft}{}{0pt}{}
\titlespacing*{\chapter}{0pt}{60pt}{36pt}
\titleformat{\section}{\normalfont\Large\bfseries\color{steel}}{\thesection}{0.8em}{}
\titleformat{\subsection}{\normalfont\large\bfseries\color{navy}}{\thesubsection}{0.8em}{}
\titleformat{\subsubsection}{\normalfont\normalsize\bfseries\color{navy}}{\thesubsubsection}{0.8em}{}
\newcommand{\usection}[1]{\section*{#1}\addcontentsline{toc}{section}{#1}}
\pagestyle{fancy}\fancyhf{}
\renewcommand{\chaptermark}[1]{\markboth{\textsc{#1}}{}}
\renewcommand{\sectionmark}[1]{\markright{#1}}
\fancyhead[L]{\small\color{navy}\nouppercase{\leftmark}}
\fancyhead[R]{\small\color{navy}\nouppercase{\rightmark}}
\fancyfoot[C]{\thepage}
\renewcommand{\headrulewidth}{0.5pt}
\renewcommand{\headrule}{{\color{navy}\hrule width\headwidth height\headrulewidth}}
\fancypagestyle{plain}{\fancyhf{}\fancyfoot[C]{\thepage}\renewcommand{\headrulewidth}{0pt}}

% ---- code et statuts epistemiques
\newcommand{\code}[1]{\texttt{\detokenize{#1}}}
\newcommand{\F}{\textcolor{ok}{\textbf{[F]}}}      % fait lu dans le code ou les mesures deposees
\newcommand{\I}{\textcolor{amber}{\textbf{[I]}}}   % inference
\newcommand{\Hy}{\textcolor{alert}{\textbf{[H]}}}  % hypothese non testee
\newcommand{\Pc}{\textcolor{ok}{\textbf{[C]}}}     % propriete justifiee par le code
\newcommand{\Pe}{\textcolor{amber}{\textbf{[E]}}}  % propriete verifiee experimentalement
\newcommand{\ie}{c.-à-d.\ }
\newcommand{\gf}{\mathbb{F}_{16}}
\newcommand{\vpshufb}{\code{vpshufb}}

% ---- theoremes
\theoremstyle{plain}
\newtheorem{proposition}{Proposition}[chapter]
\newtheorem{lemme}[proposition]{Lemme}
\newtheorem{corollaire}[proposition]{Corollaire}
\theoremstyle{definition}
\newtheorem{definition}[proposition]{Définition}
\theoremstyle{remark}
\newtheorem{remarque}[proposition]{Remarque}

\usepackage{framed}
\colorlet{shadecolor}{ice}
\newenvironment{encadre}{\begin{shaded}\small}{\end{shaded}}

\usepackage{seqsplit}
\newcommand{\sha}[1]{{\ttfamily\scriptsize\seqsplit{#1}}}

\title{Optimisation du traitement des queues de $m$-vecteurs dans les noyaux AVX2 de MAYO-C}
\author{Marcel Eric Angoah}
\date{Version provisoire du \today}

\begin{document}

% ----------------------------------------------------------------- page de titre
\begin{titlepage}
\centering
{\large École Nationale Supérieure Polytechnique de Yaoundé (ENSPY)}\\[0.3em]
{\large Master II}\\[3.5em]
{\color{muted}\rule{\linewidth}{0.6pt}}\\[1.2em]
{\Huge\bfseries\color{navy} Optimisation du traitement des queues de $m$-vecteurs dans les noyaux AVX2 de MAYO-C\par}
\vspace{0.8em}
{\Large\color{steel} L'approche A2 : un registre 128 bits pour la queue\par}
\vspace{0.8em}
{\large Analyse, validation expérimentale sur machine physique et limites}\\[1.2em]
{\color{muted}\rule{\linewidth}{0.6pt}}\\[3em]
{\large Présenté par}\\[0.3em]
{\Large\bfseries Marcel Eric Angoah}\\[2.5em]
{\large Sous la direction du}\\[0.3em]
{\Large\bfseries Dr Hervé Talé Kalachi}\\[3em]
\textbf{Version provisoire (premier jet complet)}\\[0.3em]
\today
\vfill
\begin{minipage}{0.88\linewidth}\small\color{muted}
Statut du document. Les résultats non établis figurent comme questions ouvertes ou perspectives, jamais comme contributions acquises. Chaque affirmation chiffrée renvoie à des données déposées (dépôt \code{mayo-a2-thesis}, empreintes SHA-256 en annexe~\ref{app:inventaire}).
\end{minipage}
\end{titlepage}

\pagenumbering{roman}
\chapter*{Abstract}
\addcontentsline{toc}{chapter}{Abstract}
MAYO is an Oil-and-Vinegar multivariate signature scheme submitted to post-quantum standardisation. In the official AVX2 C implementation, products of the public matrix $P_1$ by the vectors of $V$ are computed with \code{vpshufb} lookups on packed $\mathbb{F}_{16}$ vectors. When the vector length $m$ is not a multiple of 64, the official kernel covers the tail with an overlapping 32-byte load that duplicates work. We study a local change, called A2, which handles the 8-byte tail of each $P_1$ entry in a single 128-bit register and one xmm \code{vpshufb} on the low lane of the multiplication table, reducing the number of shuffles per scalar pair from 4 to 3 (MAYO-1) and from 6 to 5 (MAYO-5). We formalise the counting argument and its limits, establish the properties on which correctness rests, and measure the effect on a physical laptop with a paired protocol. Kernel gains are 12--29\,\%; the full-signature gain is 10.3\,\% for MAYO-1 but only 1.3\,\% for MAYO-5. Static analysis shows that in the signature context GCC~16.2.0 emits extra 32-byte stores in several loops, asymmetrically between the two versions, consistently with the measured loss; causality is not established. GCC's dumps show that the extra stores are per-iteration stores of the accumulator arrays that are not moved out of the loop; SRA is ruled out, and why the compiler makes this choice remains open.


\chapter*{Résumé}
\addcontentsline{toc}{chapter}{Résumé}
MAYO est un schéma de signature numérique « multivarié » de type Oil-and-Vinegar. Il est candidat à la standardisation de la cryptographie post-quantique. Dans le code C officiel (MAYO-C) pour AVX2, le programme multiplie la matrice publique $P_1$ par les vecteurs de $V$ avec l'instruction \code{vpshufb}. Les vecteurs de $\mathbb{F}_{16}$ y sont rangés à deux éléments par octet. Quand la longueur $m$ n'est pas un multiple de 64, le code officiel traite la fin du vecteur en relisant 32 octets déjà traités. Cela fait des calculs en double : 48 éléments sur 128 pour MAYO-1 ($m=80$) et 50 sur 192 pour MAYO-5 ($m=142$).

Ce mémoire étudie une petite modification, appelée A2. Elle traite la fin de 8 octets de chaque entrée de $P_1$ dans un seul registre de 128 bits : un seul \code{vpshufb} xmm remplace deux \code{vpshufb} ymm. Le nombre de \code{vpshufb} par paire de nombres passe de 4 à 3 (MAYO-1) et de 6 à 5 (MAYO-5). Nous démontrons cette limite de comptage, nous établissons les propriétés qui rendent A2 correct, puis nous mesurons l'effet sur un ordinateur portable (Intel Core i9-9880H, GCC~16.2.0). Les mesures comparent les deux versions dans les mêmes conditions, et l'unité statistique est le processus (une exécution du programme).

Voici les résultats établis. Sur les noyaux pris seuls, A2 réduit le temps de 19 à 22\,\% (\code{P1_times_Vt}) et de 12 à 29\,\% (\code{P1t_times_Vt}). Sur la signature complète, le gain est de 10,3\,\% pour MAYO-1, mais de seulement 1,3\,\% pour MAYO-5. Pour MAYO-5, il monte à 7,3\,\% quand toutes les sondes de mesure sont actives. En lisant le code produit par GCC~16.2.0 (sans l'exécuter), nous voyons que, sans toutes les sondes, le compilateur ajoute dans plusieurs boucles une écriture mémoire de 32 octets par accumulateur, et que ce n'est pas la même chose dans les deux versions. Les temps mesurés varient de la même façon, mais le lien de cause à effet n'est pas démontré, et un modèle à un seul paramètre exagère le gain d'environ 40\,\%. Ces écritures en plus sont les stockages des tableaux d'accumulateurs, refaits à chaque tour de boucle au lieu d'être faits une seule fois après la boucle. Nous écartons la passe SRA du compilateur comme cause. Pourquoi le compilateur fait ce choix reste une question ouverte.

\medskip
\noindent\textbf{Mots-clés :} MAYO, signatures multivariées, cryptographie post-quantique, AVX2, \code{vpshufb}, $\mathbb{F}_{16}$, optimisation, mesure de performance.

\tableofcontents
\listoffigures
\listoftables
\clearpage\pagenumbering{arabic}

\chapter{Introduction générale}
\label{ch:intro}

\section{Contexte}
L'algorithme de Shor~\cite{shor} rend les signatures fondées sur la factorisation et le logarithme discret vulnérables à un ordinateur quantique de grande taille. Le NIST a publié en 2024 les premières normes post-quantiques (FIPS~203, 204 et 205~\cite{nistfips}) et poursuit un appel à signatures additionnelles~\cite{nistadd}, dans lequel figurent des schémas multivariés. Parmi eux, MAYO~\cite{mayo} construit une signature de type Oil-and-Vinegar (OV)~\cite{patarin,kipnisshamir,uov} à clé publique compacte : l'application publique est une application quadratique d'un petit espace d'huile, « fouettée » (\emph{whipped}) pour produire une application plus grande. L'intérêt pratique du schéma est la petite taille des clés et des signatures ; son coût de calcul est dominé par de l'arithmétique sur $\mathbb{F}_{16}$.

Les schémas multivariés ont connu des cryptanalyses marquantes~\cite{beullensrainbow}, ce qui impose une lecture prudente de toute affirmation de sécurité ; ce mémoire ne traite pas de sécurité algébrique. Il traite de \emph{performance d'implémentation} et de \emph{propriétés d'implémentation} (correction, accès mémoire, temps constant) d'une optimisation locale du code C officiel.

\section{Parcours du travail et recentrage}
Le travail a débuté par une implémentation autonome (arithmétique $\mathbb{F}_{16}$ à tables, étude d'une fuite par canal auxiliaire de cache sur les tables, élimination de Gauss--Jordan à temps constant), documentée dans des chapitres antérieurs et motivée par les attaques temporelles de cache~\cite{kocher,bernstein}. À la demande de l'encadrant, le travail a ensuite été recentré sur l'implémentation C \emph{officielle} (MAYO-C) pour mesurer des gains qui soient comparables à l'existant. Ce mémoire porte exclusivement sur cette seconde phase. L'étude préliminaire n'est pas reprise, sauf quand un résultat négatif de cette phase éclaire une décision (chapitre~\ref{ch:discussion}).

\section{Problématique}
Dans l'implémentation AVX2 officielle, un $m$-vecteur de $\mathbb{F}_{16}$ est stocké à deux éléments par octet ($m/2$ octets) et multiplié par des scalaires à l'aide de \vpshufb{} : une instruction effectue 32 recherches dans une table de 16 octets et fournit, pour chaque élément, les produits par \emph{deux} scalaires à la fois. Un registre de 256 bits couvre donc 64 éléments. Quand $m$ n'est pas un multiple de 64, le code officiel couvre la fin du vecteur par un chargement de 32 octets reculé qui \emph{recouvre} le précédent. Ce recouvrement est documenté dans l'article d'optimisation~\cite{nibbling} pour $m=96$ (MAYO-3) : il « duplique » des opérations.

Pour les paramètres considérés ici, la duplication est importante : 48 éléments recalculés sur 128 pour MAYO-1 ($m=80$, soit 37,5\,\%) et 50 sur 192 pour MAYO-5 ($m=142$, soit 26\,\%). La question centrale est :

\begin{encadre}
\textbf{Question centrale.} La duplication de calcul due au recouvrement en fin de $m$-vecteur peut-elle être réduite par un traitement dédié de la queue, quelle en est la limite exacte en nombre de \vpshufb{}, et ce gain se traduit-il, de façon robuste et démontrable, en gain sur le temps de signature complet ?
\end{encadre}

Cette formulation sépare trois choses souvent confondues : (i) une \emph{borne de comptage} (mathématique, dans un modèle précis), (ii) sa \emph{réalisation} par un algorithme (à montrer et à tester) et (iii) l'\emph{effet sur le temps} (expérimental, dépendant du compilateur et de la machine).

\section{Questions de recherche et objectifs}
\begin{description}[leftmargin=2.2cm,labelwidth=2cm]
\item[QR1] Quel est le nombre minimal de \vpshufb{} par paire de scalaires pour un $m$-vecteur, dans la classe des traitements à tables de 16 octets, et pour quels $m$ le code officiel est-il déjà optimal ?
\item[QR2] Existe-t-il une réalisation sûre (sans lecture hors de l'entrée, sans dépendance au secret dans les adresses) de cette borne pour MAYO-1 et MAYO-5 ?
\item[QR3] Quel gain mesure-t-on sur les noyaux \code{P1_times_Vt} et \code{P1t_times_Vt}, puis sur \code{mayo_sign}, sur une machine physique, avec un protocole dont les limites sont quantifiées ?
\item[QR4] Si le gain en contexte diffère du gain isolé, que montre l'analyse du code généré, et qu'est-ce qui reste à démontrer ?
\end{description}

Les objectifs en découlent : \textbf{O1} formaliser le comptage (borne, réalisation, limite) ; \textbf{O2} documenter la conception d'A2 et les propriétés dont dépend sa correction ; \textbf{O3} mesurer avec un protocole reproductible ; \textbf{O4} caractériser l'écart entre gain isolé et gain en contexte ; \textbf{O5} situer le travail par rapport aux implémentations existantes sans revendiquer plus que ce que les sources consultées permettent.

\section{Contributions : établies et ouvertes}
Conformément à la règle de rédaction adoptée, le tableau~\ref{tab:contrib} sépare ce qui est acquis de ce qui reste ouvert.

\begin{table}[ht]\centering\small
\caption{Statut des contributions.}\label{tab:contrib}
\begin{tabularx}{\linewidth}{@{}lXl@{}}
\toprule
 & Contribution & Statut\\\midrule
C1 & Comptage : borne $\lceil m/32\rceil$, schéma officiel $2\lceil m/64\rceil$, écart $\Delta(m)$ (prop.~\ref{prop:borne}--\ref{prop:real}) & établi (preuve)\\
C2 & Réalisation A2 (queue $\le 16$ éléments) pour MAYO-1 et MAYO-5, variante SAFE sans lecture hors de P1 & établi (code, KAT, tests)\\
C3 & Gains sur noyaux isolés (12--29\,\%) et sur \code{mayo_sign} (MAYO-1 : 10,3\,\% ; MAYO-5 : 1,3\,\%) sur une machine, un compilateur & établi, avec limites\\
C4 & Localisation de la perte de gain de MAYO-5 dans \code{m_calculate} en contexte de signature & établi (mesure directe)\\
C5 & Écritures supplémentaires asymétriques dans le code généré par GCC~16.2.0, cohérentes avec C4 & fait statique établi ; causalité \emph{non} établie\\
C6 & Dépôt public de scripts, données brutes et empreintes, reproduisant les analyses (9/9) & établi\\\midrule
Q1 & Cause de ces écritures dans GCC (passe responsable, rôle de \code{temp[]}) & \emph{ouvert} (analyse E1b à exécuter)\\
Q2 & Lien causal écritures $\to$ perte de gain, avec un modèle qui n'en surestime pas l'effet & \emph{ouvert}\\
Q3 & Variante pour queues de 17 à 32 éléments ; autres compilateurs, autres microarchitectures & \emph{ouvert} (non implémenté)\\
Q4 & Temps constant au sens fort (preuve formelle) & \emph{non fourni}\\\bottomrule
\end{tabularx}
\end{table}

\section{Périmètre et exclusions}
Le travail est limité aux noyaux \code{P1_times_Vt} et \code{P1t_times_Vt} pour MAYO-1 et MAYO-5. MAYO-2 ($m=64$) et MAYO-3 ($m=108$) servent de témoins : pour eux, aucun gain n'est possible dans la classe étudiée (proposition~\ref{prop:borne}). Aucune autre optimisation n'est proposée, les noyaux optimisés ne sont pas modifiés pendant l'analyse, et la décision d'intégration dans MAYO-C est volontairement différée. L'absence d'une technique dans les sources consultées ne prouve pas qu'elle n'existe pas ailleurs : aucune nouveauté absolue n'est revendiquée.

\section{Conventions de rédaction}
Les énoncés portant sur des résultats sont étiquetés : \F{} \emph{fait} (lu dans le code généré ou dans des mesures déposées), \I{} \emph{inférence} tirée des faits, \Hy{} \emph{hypothèse} non testée. Les propriétés de sûreté du chapitre~\ref{ch:surete} utilisent \Pc{} (justifiée par le code) et \Pe{} (vérifiée expérimentalement) ; une propriété \Pe{} seule n'est pas une démonstration. Les durées sont en ticks du compteur TSC (rdtscp), non en cycles cœur.

\section{Plan}
Le chapitre~\ref{ch:etat} présente MAYO, l'arithmétique par \vpshufb{} et les implémentations existantes, et situe A2. Le chapitre~\ref{ch:a2} formalise le comptage et décrit A2. Le chapitre~\ref{ch:surete} établit les propriétés de correction et de sûreté. Le chapitre~\ref{ch:methodo} détaille la méthode expérimentale pour qu'elle soit reproductible. Le chapitre~\ref{ch:resultats} présente les résultats ; le chapitre~\ref{ch:discussion} les discute en séparant faits, inférences et hypothèses. La conclusion récapitule contributions, limites et perspectives.

\chapter{Contexte, état de l'art et positionnement}
\label{ch:etat}

\section{Le schéma MAYO et ses noyaux de calcul}
\subsection{Structure}
MAYO~\cite{mayo} dérive d'un schéma UOV~\cite{uov} dont l'espace d'huile est volontairement petit, ce qui réduit la clé publique ; la sécurité est restaurée en composant $k$ copies de l'application publique (« fouettage »). Les paramètres de ce mémoire sont ceux de l'arbre officiel de référence (annexe~\ref{app:inventaire}). Les grandeurs utiles à l'étude des noyaux sont données au tableau~\ref{tab:params}.

\begin{table}[ht]\centering\small
\caption{Paramètres des noyaux étudiés. $m$ : longueur du $m$-vecteur ; $b=m/2$ : octets par $m$-vecteur ; $L$ : limbes de 64 bits d'un accumulateur ; $K$ : nombre de vecteurs de $V$ ; $V$ : dimension de la partie vinaigre ; $N_{\mathrm{ent}}=V(V+1)/2$ : nombre d'entrées de $P_1$.}\label{tab:params}
\begin{tabular}{@{}lrrrrrrr@{}}\toprule
Paramètre & $m$ & $b$ & $L$ & $K$ & $V$ & $N_{\mathrm{ent}}$ & $m \bmod 64$\\\midrule
MAYO-1 & 80 & 40 & 5 & 10 & 80 & 3\,240 & 16\\
MAYO-2 & 64 & 32 & 4 & -- & 73 & -- & 0\\
MAYO-3 & 108 & 54 & -- & -- & -- & -- & 44\\
MAYO-5 & 142 & 71 & 9 & 12 & 142 & 10\,153 & 14\\\bottomrule
\end{tabular}
\end{table}

\subsection{Rôle de $P_1$}
$P_1$ est une matrice triangulaire supérieure de $m$-vecteurs de $\mathbb{F}_{16}^m$ ; elle est stockée à plat, une entrée après l'autre, à la position $e\cdot b$ pour l'entrée numéro $e$ (disposition empaquetée du flux de clé). Les noyaux \code{P1_times_Vt} et \code{P1t_times_Vt} accumulent, pour chaque entrée de $P_1$ et chaque vecteur de $V$, le produit d'un scalaire de $\mathbb{F}_{16}$ par le $m$-vecteur de l'entrée, dans un accumulateur de $L$ limbes ; le premier traite $P_1\cdot V^{t}$, le second la contribution transposée. Leurs bornes de boucle ($V$, $K$) et les adresses de chargement ($e\cdot b$) ne dépendent que de paramètres publics. Ces deux noyaux sont appelés par \code{m_calculate_M_and_VPV}, elle-même appelée par la signature (chapitre~\ref{ch:resultats} pour leur poids).

\section{Arithmétique $\mathbb{F}_{16}$ par \vpshufb}
\subsection{Principe}
Avec des éléments de 4 bits, la multiplication d'un vecteur par un scalaire $c$ est une recherche dans une table de 16 entrées $x\mapsto c\cdot x$. L'instruction \vpshufb{} (\code{_mm256_shuffle_epi8}) réalise, pour chaque octet d'un registre d'index, une recherche dans une table de 16 octets (celle de la voie de 128 bits qui contient l'octet) ; si le bit 7 de l'index est nul, l'indice est le quartet bas, sinon la sortie est 0. C'est la méthode des « quatre Russes »~\cite{russes} à $t=1$ (voir aussi~\cite{albrecht2011} pour $\mathbb{F}_2$).

\subsection{Deux produits par recherche}
Pour une paire de scalaires $(c_{2p},c_{2p+1})$, la table $T_p[x]=(c_{2p}x)\oplus\big((c_{2p+1}x)\ll4\big)$ fournit \emph{les deux} produits dans un seul octet. Le code stocke $T_p$ dans un registre 256 bits dont les deux voies sont égales (lemme~\ref{lem:voies}). Un registre de données de 32 octets contient 64 éléments ; il faut séparer les quartets impairs et pairs (ET avec \code{0x0f}, décalage de 4) et donc deux \vpshufb{} par registre et par paire. Les accumulateurs impairs et pairs sont maintenus séparés et ne sont « désintriqués » qu'à la fin de la ligne, technique décrite dans l'article \emph{Nibbling MAYO}~\cite{nibbling}.

\subsection{Recouvrement en fin de vecteur}
Pour $m\le64$, $64<m\le128$ et $128<m\le192$, le code officiel utilise respectivement 1, 2 et 3 chargements de 32 octets, le dernier étant \emph{reculé} pour finir à l'octet $b$ (adresse \code{vec + M_VEC_LIMBS_MAX - 4} limbes) ; les octets en commun sont donc traités deux fois et les stockages finaux se recouvrent avec des valeurs identiques. Le décompte pour MAYO-5 est le suivant : chargements $[0,32)$, $[32,64)$ et $[40,72)$, soit 192 éléments calculés pour 142 utiles ; le dernier chargement lit un octet au-delà de l'entrée de $P_1$ (figure~\ref{fig:charge}).

\begin{figure}[ht]\centering
\begin{tikzpicture}[x=0.105cm,y=1cm,>=Stealth,font=\small]
 \def\h{0.45}
 % echelle : 1 unite = 1 octet, entree de 71 octets : 0..71
 \node[anchor=east] at (-2,4.2) {Entrée de $P_1$ (71 octets)};
 \fill[ice] (0,4) rectangle (71,4.5); \draw (0,4) rectangle (71,4.5);
 \foreach \x in {0,32,64,71}{\draw[muted] (\x,3.9)--(\x,4.6); \node[font=\scriptsize,muted] at (\x,4.85) {\x};}
 % officiel
 \node[anchor=east] at (-2,3) {Officiel : 3 ymm};
 \fill[steel!35] (0,2.8) rectangle (32,3.25); \draw (0,2.8) rectangle (32,3.25);
 \fill[steel!55] (32,2.8) rectangle (64,3.25); \draw (32,2.8) rectangle (64,3.25);
 \fill[alert!40] (40,2.25) rectangle (72,2.7); \draw (40,2.25) rectangle (72,2.7);
 \node[font=\scriptsize] at (16,3.02) {$[0,32)$}; \node[font=\scriptsize] at (48,3.02) {$[32,64)$}; \node[font=\scriptsize] at (56,2.47) {$[40,72)$ : recouvre $[40,64)$, lit 1 octet de trop};
 % A2
 \node[anchor=east] at (-2,1.2) {A2 SAFE : 2 ymm + 1 xmm};
 \fill[steel!35] (0,1) rectangle (32,1.45); \draw (0,1) rectangle (32,1.45);
 \fill[steel!55] (32,1) rectangle (64,1.45); \draw (32,1) rectangle (64,1.45);
 \fill[ok!40] (63,0.45) rectangle (71,0.9); \draw (63,0.45) rectangle (71,0.9);
 \node[font=\scriptsize] at (16,1.22) {$[0,32)$}; \node[font=\scriptsize] at (48,1.22) {$[32,64)$};
 \node[font=\scriptsize,anchor=east] at (62,0.67) {queue $[63,71)$, puis décalage de 8 bits};
\end{tikzpicture}
\caption{Chargements pour une entrée de $P_1$ de MAYO-5 (octets). Le noyau officiel recouvre l'intervalle $[40,64)$ et dépasse l'entrée d'un octet ; A2 SAFE lit exactement $[0,71)$.}\label{fig:charge}
\end{figure}

\section{Implémentations existantes}
\begin{description}[leftmargin=0.5cm]
\item[Référence C.] Arithmétique scalaire, sans \vpshufb{}.
\item[AVX2 officiel (MAYO-C).] Fichier \code{src/AVX2/shuffle_arithmetic.h} : macros \code{MULTIPLY_ACCUMULATE_M_VEC} et \code{DISENTANGLE_AND_ACCUMULATE_M_VECS}, recouvrement décrit plus haut. Base de ce travail : commit amont \code{74de637}~\cite{mayoc}.
\item[\emph{Nibbling MAYO}~\cite{nibbling}.] Implémentation AVX2 et Cortex-M4 optimisée avec représentation en quartets ; §5.1 : un \vpshufb{} fournit 64 multiplications, accumulation séparée impair/pair, et, pour $m=96$ (MAYO-3 de l'époque), un recouvrement qui « duplique 32 opérations » ; débit d'un \vpshufb{} par cycle sur Skylake, deux sur Ice Lake. Le code artefact (\code{shuffle_arithmetic_\{64,96,128\}.h}) charge à des positions de limbes 0 et 2 (octets $[0,32)$ et $[16,48)$) et utilise quatre \vpshufb{} par paire pour $m=96$.
\item[Backend GFNI.] Utilise \code{vgf2p8affineqb} et lit en pleine largeur dans une marge (\code{MAYO_MVEC_SLACK}) ; hors du cadre de ce mémoire.
\end{description}

\section{Positionnement d'A2}
Le tableau~\ref{tab:etat} sépare ce qui était connu, ce qui est postérieur à l'article et ce qu'A2 fait différemment, dans les sources consultées (article §5.1 lu en entier et les parties AVX2 du code ; historique complet de \code{shuffle_arithmetic.h} dans le dépôt officiel, commits \code{3c9be9e}, \code{d81b689}, \code{11ff0f0}, \code{4b7cd94}, \code{fef05b5} et tête de \code{main} \code{74de637}).

\begin{table}[ht]\centering\small
\caption{Situation d'A2 par rapport aux sources consultées.}\label{tab:etat}
\begin{tabularx}{\linewidth}{@{}lXX@{}}\toprule
Catégorie & Élément & Source\\\midrule
Connu avant A2 & \vpshufb{} à $t=1$, deux produits par octet ; séparation impair/pair ; désintrication tardive & \cite{nibbling}\\
Connu avant A2 & Recouvrement de registres pour $m$ non multiple de 64 (MAYO-3, $m=96$) & \cite{nibbling}\\
Postérieur à l'article & Généralisation par chargements reculés pour tout $m\le192$ & code officiel (Round~2)\\
Différent dans A2 & Queue $\le16$ éléments regroupée dans \emph{un} xmm ([bas$\mid$haut] des quartets), \emph{un} \vpshufb{} xmm sur la voie basse de la table & ce travail\\
Différent dans A2 & Variante SAFE pour $m=142$ : charge $[63,71)$ puis décale, aucune lecture hors de l'entrée & ce travail\\
Différent dans A2 & Propriétés vérifiées de la queue (chapitre~\ref{ch:surete}) & ce travail\\
Non revendiqué & Recouvrement, séparation impair/pair, désintrication tardive, schéma $t=1$, nouveauté absolue & --\\\bottomrule
\end{tabularx}
\end{table}

Aucun \code{vpshufb} xmm sur une queue n'a été trouvé dans l'artefact de l'article, ni dans l'historique officiel jusqu'à la tête de \code{main} : les \code{__m128i} qui apparaissent à partir de \code{fef05b5} (lignes 548--606) concernent la transposition par décalage de bits de $m$-vecs et le chargement d'une table de 16 octets diffusée sur 256 bits. \textbf{Limite :} l'article n'a été lu intégralement que dans sa section 5.1 et ses parties AVX2 pertinentes ; l'absence de la technique dans ces sources ne prouve pas qu'elle n'existe pas ailleurs.

\section{Canaux auxiliaires et temps constant}
Les attaques temporelles~\cite{kocher} et de cache~\cite{bernstein} imposent que les adresses et les branches ne dépendent pas de données secrètes. Dans les noyaux, les tables de multiplication dérivées du vinaigre ne sont que des opérandes de données de \vpshufb{} : ni adresses, ni conditions. A2 doit conserver cette propriété ; l'état de la preuve est discuté au chapitre~\ref{ch:surete}.

\chapter{L'optimisation A2 : comptage, conception, hypothèses et limites}
\label{ch:a2}

Ce chapitre répond à QR1 et décrit la réalisation (QR2). Les propositions~\ref{prop:borne} à~\ref{prop:real} sont des énoncés mathématiques \emph{dans un modèle}; aucune conclusion de performance n'en est déduite (section~\ref{sec:perfborne}).

\section{Modèle de comptage}
\begin{definition}[Données]
Un \emph{$m$-vec} est un vecteur $v\in\mathbb{F}_{16}^m$, $m$ pair, stocké sur $b=m/2$ octets : $v_{2i}$ est le quartet bas de l'octet $i$, $v_{2i+1}$ le quartet haut. Les $K$ scalaires $c_0,\dots,c_{K-1}$ sont regroupés en $P=\lceil K/2\rceil$ paires (complétées par 0 si $K$ est impair). Pour chaque $k<K$ on calcule $\mathrm{acc}_k\leftarrow\mathrm{acc}_k\oplus c_k v$.
\end{definition}
Pour la paire $p$, $T_p[x]=(c_{2p}x)\oplus((c_{2p+1}x)\ll4)$ ; \vpshufb{} de largeur $W\in\{16,32\}$ octets effectue $W$ recherches indépendantes, chacune livrant, pour \emph{un} élément de $v$, les deux produits de la paire.

\begin{definition}[Classe $\mathcal{C}$]
Algorithmes qui (i) obtiennent chaque produit $c_kv_j$ par une recherche dans une table de 16 octets ; (ii) fournissent à chaque recherche un quartet isolé dans un octet (bit~7 nul) ; (iii) ne font pas deux recherches pour le même couple (élément, paire) hors du recouvrement entre registres ; (iv) utilisent librement XOR, ET, décalages, permutations, chargements. Sont \emph{hors} de $\mathcal{C}$ : GFNI, représentations \emph{bitsliced}, tables de 256 octets.
\end{definition}
La quantité comptée est $N(\text{algo},m)$ : nombre de \vpshufb{} exécutés \emph{par paire de scalaires et par $m$-vec}.

\section{Énoncés}
\begin{proposition}[Borne]\label{prop:borne}
Tout algorithme de $\mathcal{C}$ vérifie $N\ge\lceil m/32\rceil$.
\end{proposition}
\begin{proof}
Chaque élément $v_j$ ($j<m$) exige au moins une recherche par paire ; une instruction en fournit au plus $W\le32$. Donc $32N\ge m$.
\end{proof}

\begin{proposition}[Schéma officiel avec recouvrement]\label{prop:off}
Le schéma qui couvre les $b$ octets par $\lceil b/32\rceil$ chargements de 32 octets, le dernier reculé, et traite chaque registre par deux \vpshufb{} (impairs, pairs) a $N_{\mathrm{off}}(m)=2\lceil m/64\rceil$.
\end{proposition}
\begin{proof}
$\lceil b/32\rceil=\lceil m/64\rceil$ chargements, chacun fournissant 32 indices impairs et 32 pairs, soit deux \vpshufb{}.
\end{proof}

\begin{proposition}[Réalisation de la borne si $0<r\le32$]\label{prop:real}
Écrivons $m=64q+r$, $0\le r<64$, et $t=r/2$ le nombre d'octets de la queue. Si $0<r\le32$, l'algorithme suivant est dans $\mathcal{C}$ et a $N=2q+1=\lceil m/32\rceil$ :
\begin{enumerate}[nosep]
\item pour $g<q$ : charger l'intervalle $[32g,32g+32)$ et le traiter par deux \vpshufb{} ymm ;
\item pour la queue : charger les $t$ octets de $[32q,32q+t)$, construire $X=[\text{quartets bas des }t\text{ octets}\mid\text{quartets hauts des }t\text{ octets}]$ ($r\le32$ indices dans un seul registre) et appliquer \emph{un} \vpshufb{} ;
\item réassembler les sorties de la queue (même formule qu'un registre complet, restreinte aux $r$ positions).
\end{enumerate}
\end{proposition}
\begin{proof}
Le décompte est $2q+1$, et $\lceil m/32\rceil=2q+\lceil r/32\rceil=2q+1$ pour $0<r\le32$.
\end{proof}
\emph{Conditions de réalisation} (à satisfaire par le code, non par le modèle) : (a) le chargement de la queue ne lit que des octets de l'entrée ; (b) si $r\le16$ ($t\le8$), $X$ tient dans un registre 128 bits et la table utilisée est la \emph{voie basse} de $T_p$ ; si $16<r\le32$, $X$ occupe un registre 256 bits (voie basse = impairs, haute = pairs), ce qui exige l'égalité des deux voies ; (c) le surcoût de préparation de la queue est payé une fois par entrée, celui du réassemblage une fois par paire.

\begin{corollaire}[Limite exacte]
Pour $m\le192$, $\Delta(m):=N_{\mathrm{off}}(m)-\lceil m/32\rceil=1$ si $0<(m\bmod64)\le32$ et $0$ sinon. Dans le second cas, le code officiel atteint déjà la borne : aucun algorithme de $\mathcal{C}$ ne peut faire mieux en nombre de \vpshufb{}.
\end{corollaire}

\begin{table}[ht]\centering\small
\caption{Application aux paramètres actuels.}\label{tab:delta}
\begin{tabular}{@{}lrrrrrrl@{}}\toprule
 & $m$ & $q$ & $r$ & $N_{\mathrm{off}}$ & $\lceil m/32\rceil$ & $\Delta$ & Statut\\\midrule
MAYO-1 & 80 & 1 & 16 & 4 & 3 & 1 & implémentée (A2), queue de 8 octets\\
MAYO-2 & 64 & 1 & 0 & 2 & 2 & 0 & sans objet\\
MAYO-3 & 108 & 1 & 44 & 4 & 4 & 0 & impossible dans $\mathcal{C}$\\
MAYO-5 & 142 & 2 & 14 & 6 & 5 & 1 & implémentée (A2), queue de 7 octets + 1 de bourrage\\\bottomrule
\end{tabular}
\end{table}

\paragraph{Ce qui est implémenté et ce qui ne l'est pas.} Le prototype traite les queues d'au plus 16 éléments ($t\le8$, $r\le16$). La variante pour $16<r\le32$ \textbf{n'est ni implémentée ni testée} ; elle n'est pas un résultat acquis. Aucun des quatre paramètres actuels n'a $r\in(16,32]$.

\section{Conception d'A2}
\subsection{Structure du noyau}
Pour une entrée de $P_1$ et une ligne d'accumulateurs de $K$ vecteurs, A2 conserve inchangés les $q=\lfloor m/64\rfloor$ registres ymm complets sur les $32q$ premiers octets, puis traite la queue ainsi :
\begin{enumerate}[nosep]
\item \textbf{Préparation (une fois par entrée)} : chargement de 8 octets (\code{_mm_loadl_epi64}), décalage de 4 bits (\code{vpsrlw}), entrelacement \code{vpunpcklqdq} (quartets bas dans les octets 0--7, hauts dans 8--15), ET avec \code{0x0f} : quatre instructions (cinq dans la variante SAFE, qui ajoute un décalage d'octets).
\item \textbf{Produits} : un \code{vpshufb} xmm avec la \emph{moitié basse} de la table de multiplication déjà chargée, et un \code{vpxor} xmm dans un accumulateur xmm.
\item \textbf{Réassemblage (fin de ligne)} : séparation impair/pair identique à la partie ymm, et lecture-modification-écriture d'\emph{un} limbe de 64 bits (le limbe d'indice $4q=L-1$, soit 4 pour MAYO-1 et 8 pour MAYO-5).
\end{enumerate}
Dans l'assembleur de GCC~16.2.0, la boucle interne passe donc, par entrée, de 20 à 15 \code{vpshufb} pour MAYO-1 (10 ymm + 5 xmm) et de 36 à 30 pour MAYO-5 (24 ymm + 6 xmm).

\subsection{Variante par défaut et variante SAFE (MAYO-5)}
Pour MAYO-1, $b=40=32+8$ : la queue est exactement la fin de l'entrée, $[32,40)$. Pour MAYO-5, $b=71$ et $64+8=72$ :
\begin{itemize}[nosep]
\item \emph{défaut} : lit $[64,72)$, soit \textbf{un octet au-delà} de l'entrée (comme le noyau officiel, dont le dernier chargement couvre $[40,72)$) ;
\item \emph{SAFE} : lit $[63,71)$ puis décale le registre de 8 bits vers les poids faibles ; l'octet 63 est éliminé, les octets 64--70 sont placés aux positions 0--6, la position 7 reçoit 0. \textbf{Toute lecture reste dans l'entrée.}
\end{itemize}
La variante retenue pour une éventuelle intégration est SAFE seule. Les quartets de bourrage (positions 142 et 143 de MAYO-5, car $16L=144$) n'ont pas d'influence sur les positions utiles (section~\ref{sec:bourrage}).

\subsection{Préconditions}
(P1) $P_1$ est dans la disposition empaquetée, sans bourrage entre entrées, à l'adresse $P_1+e\cdot b$ ; (P2) les accumulateurs ont $L$ limbes avec bourrage de quartets seulement en fin de dernier limbe ; (P3) les tables de multiplication sont construites par la routine officielle (voies égales, lemme~\ref{lem:voies}) ; (P4) $m\in\{80,142\}$ : les noyaux A2 ne sont pas actifs pour les autres paramètres (condition de compilation \code{MAYO_TAIL_XMM_ACTIVE}) ; (P5) compilation sans \code{ENABLE_PARAMS_DYNAMIC}.

\section{La borne n'implique pas l'accélération}\label{sec:perfborne}
Une réduction du nombre de \vpshufb{} ne garantit pas une réduction du temps : s'y ajoutent le coût de préparation de la queue, les dépendances, les rechargements depuis la pile, les chargements non alignés et le débit des unités. Si seule la contrainte de débit du \vpshufb{} (un par cycle sur Skylake~\cite{nibbling}) bornait la boucle, le rapport de temps attendu serait $N_{A2}/N_{\mathrm{off}}$ : $15/20=0{,}750$ (MAYO-1) et $30/36=0{,}833$ (MAYO-5).

\begin{table}[ht]\centering\small
\caption{Rapport attendu par le seul débit de \vpshufb{} et rapport mesuré (S1, A2/officiel, noyaux isolés).}\label{tab:borne-mesure}
\begin{tabular}{@{}lccl@{}}\toprule
Noyau & Borne « débit seul » & Mesuré (S1) & Lecture\\\midrule
MAYO-1 \code{P1_times_Vt} & 0,750 & 0,780 & plus lent que la borne\\
MAYO-1 \code{P1t_times_Vt} & 0,750 & 0,710 & plus rapide que la borne\\
MAYO-5 \code{P1_times_Vt} & 0,833 & 0,807 & plus rapide que la borne\\
MAYO-5 \code{P1t_times_Vt} & 0,833 & 0,880 & plus lent que la borne\\\bottomrule
\end{tabular}
\end{table}
Deux cas sur quatre sont plus rapides que ne le permet le seul comptage de \vpshufb{} : d'autres mécanismes contribuent (par exemple moins de rechargements de pile : 11 $\to$ 5 pour MAYO-1 \code{P1_times_Vt}, 30 $\to$ 24 pour MAYO-5), sans que cela soit établi. Le modèle de débit n'explique donc pas à lui seul les gains. Une hypothèse de performance H-P (« sur une microarchitecture où le débit des shuffles est le goulot, le temps est borné inférieurement par $N\cdot P$ cycles par entrée ») ne peut se tester qu'avec des cycles cœur, pas des ticks TSC ; elle est postérieure à ce travail et reste ouverte.

\chapter{Correction et propriétés de sûreté}
\label{ch:surete}

Ce chapitre répond à la partie « sûreté » de QR2. Chaque propriété porte son étiquette : \Pc{} (déduite du code et de ses invariants, par relecture et assertions de compilation), \Pe{} (vérifiée expérimentalement) ou \Hy{} (hypothèse sur la machine ou le compilateur). Une propriété \Pe{} seule n'est pas une démonstration. Notation : $n_y=\lfloor m/64\rfloor$ (1 pour MAYO-1, 2 pour MAYO-5), $T_{\mathrm{OFF}}=32n_y$ (32 ; 64), $T_{\mathrm{LIMB}}=4n_y$ (4 ; 8).

\section{Bornes de lecture de $P_1$}
L'entrée $e$ occupe les octets $[eb,(e+1)b)$.
\begin{itemize}
\item \Pc{} Les $n_y$ chargements ymm lisent $[eb+32g,eb+32g+32)$, $g<n_y$ ; comme $32n_y\le b$ ($32\le40$ ; $64\le71$) ils restent dans l'entrée.
\item \Pc{} MAYO-1 : la queue lit $[eb+32,eb+40)$, exactement la fin de l'entrée.
\item \Pc{} MAYO-5, SAFE : la queue lit $[eb+63,eb+71)$ puis décale de 8 bits ; toute lecture reste dans l'entrée.
\item \Pc{} (variante par défaut, non retenue) MAYO-5 lit $[eb+64,eb+72)$ : un octet au-delà. Pour $e<N_{\mathrm{ent}}-1$, c'est le premier octet de l'entrée suivante ; pour la dernière, c'est l'octet 720\,863 du tampon, alloué (zone de $8\cdot$\code{P1_LIMBS_MAX}$=731\,016$ octets). Des \code{_Static_assert} imposent $N_{\mathrm{ent}}b+\text{OVERREAD}\le8\cdot$\code{P1_LIMBS_MAX} et OVERREAD$\le1$.
\item \Pe{} Test « garde-page » (fin de $P_1$ collée à une page interdite, marge 0) : l'officiel et A2 par défaut déclenchent SIGSEGV en MAYO-5 ; A2 SAFE non ; MAYO-1 : aucun des trois. Ce test contrôle la propriété ci-dessus sans la démontrer.
\end{itemize}
Le noyau officiel de MAYO-5 lit lui aussi 1 octet au-delà (dernier ymm sur les octets 40--71) et celui de MAYO-3 en lit 2 (second ymm sur $[24,56)$ pour $b=54$). A2 SAFE ne dépend d'aucune marge.

\section{Bornes d'écriture}
Pour la ligne $r$, l'accumulateur est à $\mathrm{acc}+rKL$ (limbes) ; le vecteur $k$ occupe $[kL,(k+1)L)$.
\begin{itemize}
\item \Pc{} Les stockages ymm écrivent les limbes $[kL+4g,kL+4g+4)$, $g<n_y$, jamais au-delà du limbe $4n_y-1\le L-1$.
\item \Pc{} La queue lit-modifie-écrit un seul limbe de 64 bits, d'indice $T_{\mathrm{LIMB}}=L-1$, pour $k<K$ (et $k+1<K$ pour le second, garde explicite). Toutes les écritures sont donc dans les $VKL$ limbes de \code{acc}. Le noyau officiel écrit un ymm à $L-4$ : même étendue.
\item \Pe{} Un « canari » placé après \code{acc} n'est pas modifié ; $P_1$, passé en \code{const}, n'est pas modifié (copie de l'entrée inchangée).
\end{itemize}

\section{Invariant des deux voies des tables}
A2 lit la \emph{voie basse} de chaque table et suppose qu'elle égale la voie haute.
\begin{lemme}\label{lem:voies}\Pc{} Pour tout $i$, \code{V_multabs[i]} a ses deux voies de 128 bits égales.
\end{lemme}
\emph{Justification.} (i) Chacune des quatre lignes de 32 octets de \code{__gf16_mulbase} répète deux fois la même suite de 16 octets (constantes lues dans \code{arithmetic_common.h}, lignes 12--17) ; (ii) \code{tbl32_gf16_multab2(b)} est un OU exclusif de ET entre ces lignes et des masques obtenus par \code{_mm256_set1_epi16(b & 0xf)} décalé et comparé, opérations agissant par élément de 16 bits donc par voie, avec des masques identiques dans les deux voies ; (iii) \code{mayo_V_multabs} combine deux telles tables par \code{^} et \code{_mm256_slli_epi16(.,4)}, également par voie. Par induction l'égalité est conservée.

\Pe{} 272 tables vérifiées par le mode \code{MAYO_TAIL_XMM_CHECK} (arrêt en cas d'écart) sur les tables passées aux noyaux. C'est une \emph{précondition} d'A2 : si les tables étaient construites autrement (autre routine, autre backend), l'invariant devrait être revérifié.

\section{Influence des quartets de bourrage}\label{sec:bourrage}
Seul MAYO-5 a des quartets de bourrage (positions 142 et 143 ; MAYO-1 : $80=5\cdot16$).
\begin{itemize}
\item \Pc{} \emph{Indépendance positionnelle.} La sortie à la position $j$ ne dépend que des entrées à la position $j$ de chaque $m$-vec : une recherche lit un index et une table, jamais un autre index ; les décalages \code{srli/slli_epi16} mélangent deux octets voisins mais chaque résultat est suivi d'un ET avec \code{0x0f} (ou provient d'un registre déjà masqué).
\item \Pc{} Les deux seules différences entre officiel et A2 aux positions de bourrage sont : officiel : produits de l'octet 71 (premier octet de l'entrée suivante) ; A2 SAFE : 0.
\item \Pc{} Les noyaux alimentent \code{Pv}, puis \code{Vt_times_Pv} produit \code{VP1V}, consommé par \code{compute_rhs}, premier consommateur qui mélange des positions ; il annule les quartets de bourrage avant usage (\code{src/mayo.c}, lignes 50--60, masque $(1\ll((m\bmod16)\cdot4))-1$). Entre les noyaux et \code{compute_rhs}, les opérations sont positionnelles.
\item \Pe{} Tests unitaires (20\,000 cas par taille, 8 familles d'entrées) : 0 écart sur les $M/2$ octets utiles et sur les limbes complets après masquage ; sur 30\,010 signatures appariées de MAYO-5 (15\,005 pour MAYO-1 et MAYO-2), 0 signature différente.
\end{itemize}
\textbf{Limite.} Cette propriété suppose que \code{Pv} n'est consommé que par les chemins énumérés, par recherche dans le source du backend AVX2 ; une modification future qui lirait le bourrage de \code{Pv} devrait être réexaminée.

\section{Temps constant : ce qui est établi et ce qui ne l'est pas}
\begin{itemize}
\item \Pc{} Bornes de boucle ($V$, $K$) et adresses de chargement ($eb$) ne dépendent que de paramètres publics ; les tables dérivées du vinaigre sont des opérandes de données de \vpshufb{}, jamais des adresses ni des conditions. L'officiel a la même structure.
\item \Pc{}$\to$à vérifier sur le binaire : A2 introduit, au réassemblage de la queue, deux \code{vmovq} d'un registre vectoriel vers un registre général (\code{ar[...] ^= (uint64_t)_mm_cvtsi128_si64(...)}) : des données dérivées du secret transitent par un registre général, uniquement comme opérande d'un XOR vers la mémoire. L'inspection de l'assembleur (script \code{inspect_asm.py}) ne trouve aucun usage non conforme (adresse, branche, drapeau) : résultat pour \emph{une} version de GCC, pas une garantie générale.
\item \Pe{} (indicateur) Valgrind en mode temps constant (données secrètes marquées indéfinies), dans la VM seulement ; il détecte branches et adresses dépendant de données marquées sur les exécutions testées.
\item \Hy{} Les instructions \vpshufb{}, \code{vpxor}, \code{vpand}, décalages, chargements et stockages ont un temps indépendant des données sur les microarchitectures visées (documentation constructeur, non vérifiée ici).
\end{itemize}
\textbf{Aucune preuve formelle de temps constant n'est fournie.} La réussite des KAT et l'absence de différence fonctionnelle ne démontrent pas le temps constant.

\section{Correction fonctionnelle}
KAT officiels réussis pour MAYO-1, 2, 3 et 5 (toutes les sessions), tests de schéma réussis, tests unitaires précédents, zéro signature différente dans toutes les paires A/B, nombre de tentatives de signature identique. La correction est donc établie \emph{sur les entrées testées} ; l'argument général repose sur les propriétés \Pc{} ci-dessus.

\chapter{Méthodologie expérimentale}
\label{ch:methodo}

L'objectif est qu'un lecteur disposant du dépôt, des archives et d'une machine comparable puisse refaire chaque analyse et, à des variations de bruit près, chaque mesure. Les commandes figurent à l'annexe~\ref{app:inventaire}.

\section{Plateforme}
\paragraph{Machine physique (toutes les mesures citées au chapitre~\ref{ch:resultats}).} Portable Intel Core i9-9880H (8 cœurs, 16 threads, microcode \code{0x104}) ; \code{systemd-detect-virt} = \code{none}, drapeau hyperviseur = 0 ; Kali 7.1.5 (démarrage sur disque, double amorçage) ; GCC 16.2.0 (Debian 16.2.0-3), cmake 4.3.4, Python 3.14.8, numpy 2.4.6. Gouverneur \code{performance}, turbo actif, SMT actif, secteur branché. Processus épinglé sur le cœur 1 (\code{taskset}). La fréquence instantanée du cœur n'est enregistrée qu'au lancement (de 800 MHz au repos à 4\,540 MHz selon les sessions) ; elle n'est pas une fréquence soutenue. Horloge du portable à UTC+2 ; les journaux \code{machine.txt} sont en UTC.

\paragraph{Environnement virtuel (historique, non utilisé pour les résultats de performance).} VM KVM à 2 cœurs, GCC 13.3.0, cmake 3.28.3. Les mesures de cette VM (qui ont orienté les décisions de conception) ne sont jamais présentées comme performances matérielles.

\paragraph{Compilation.} Options de la bibliothèque : \code{-O3 -funroll-loops -march=haswell -maes}, identiques pour les deux versions comparées. Pour l'analyse statique, la commande exacte de chaque objet est lue dans \code{compile_commands.json}.

\section{Code comparé}
Base amont : dépôt officiel \code{PQCMayo/MAYO-C}, commit \code{74de637a7c4a7f9fafa888183bc05f5cf920cd2d}. Les modifications sont appliquées par \code{git am} dans l'ordre : \code{0001} (prototype A2), \code{0002} (sondes de décomposition, inertes par défaut), \code{0003} (contrôle de l'effet des sondes), \code{0004} (analyse statique E1, aucun noyau modifié). Les identifiants de commit changent avec \code{git am} ; les arbres, non (annexe~\ref{app:inventaire}). L'arbre des mesures S1/S2 est celui après \code{0001} ; celui de PH1/CT1 après \code{0003}. Les noyaux optimisés n'ont pas été modifiés pendant la phase d'analyse.

\section{Protocole de mesure apparié}
Pour chaque niveau mesuré (\code{P1_times_Vt}, \code{P1t_times_Vt}, \code{m_calculate}, \code{mayo_sign}), un processus exécute des paires (officiel, A2) en ordre \emph{aléatoire}, avec la même graine et le même message dans les deux bras. Le temps est lu par \code{rdtscp} (ticks TSC) ; l'ordre d'exécution de chaque paire est enregistré (\code{xmm_first}). Un processus comprend 3\,001 paires. Les signatures des deux bras sont comparées octet à octet et le nombre de tentatives de signature est vérifié.

\section{Unité statistique et intervalles}
Les 3\,001 paires d'un même processus ne sont pas des répétitions indépendantes (même état de la machine, même fréquence) : l'IC95 « naïf » calculé sur les paires est de l'ordre de $\pm0{,}0001$ pour \code{mayo_sign}, alors que la dispersion entre processus est 5 à 15 fois plus grande. \textbf{L'unité statistique est le processus.} Pour chaque niveau et chaque processus on calcule la médiane du rapport A2/officiel ; la valeur rapportée est la médiane de ces médianes, avec un IC95 \emph{bootstrap sur les processus}. Une seule différence entre sessions est testée (S2$-$S1) : avec deux démarrages, la composante « entre sessions » n'est pas estimable séparément ; on donne l'écart et son IC bootstrap.

\paragraph{Effet d'ordre.} Pour chaque processus : médiane du rapport sur les paires « A2 en premier » et sur les paires « officiel en premier » ; l'effet d'ordre est la différence, avec IC bootstrap sur les processus ; il est jugé relativement au gain.

\section{Sessions}
\begin{table}[ht]\centering\small
\caption{Sessions de mesure sur machine physique (dates UTC).}\label{tab:sessions}
\begin{tabularx}{\linewidth}{@{}lXXX@{}}\toprule
Session & Date, démarrage & Plan & Objet\\\midrule
S1 & 19 oct. 21:58--22:43 ; un démarrage & 4 campagnes $\times$ 5 processus $\times$ 3\,001 paires (graines 1--20) & gains\\
S2 & 20 oct. 12:10 ; nouveau démarrage & 5 processus (graines 21--25) & reproductibilité\\
PH1 & 21 oct. 19:13 ; nouveau démarrage & MAYO-5 : 10 processus ; MAYO-1, MAYO-2 : 5 & décomposition instrumentée\\
CT1 & 21 oct. 19:28 ; même démarrage que PH1 & 5 configurations $\times$ 10 processus (graines 41--50) & effet des sondes\\\bottomrule
\end{tabularx}
\end{table}
Les scripts de lancement refusent d'écraser un dossier ou une archive existants ; aucune campagne n'écrase la précédente. S1 et S2 comparent deux démarrages distincts avec le même code, la même machine et le même compilateur. Temps absolus : S2 est environ 7\,\% plus lent que S1 pour tous les niveaux, témoins compris ; la cause n'est pas mesurée et les rapports ne sont pas affectés.

\section{Instrumentation de décomposition}
Des sondes (\code{rdtscp} précédé d'un \code{lfence}) mesurent, pour chaque signature et chaque bras, le temps total, le nombre de tentatives et, pour chaque phase, le temps et le nombre d'appels. Elles sont inertes par défaut : sans la macro correspondante, le code objet de \code{mayo.c} et \code{arithmetic.c} est identique à l'original (vérifié par \code{cmp} et \code{objdump}). Le patch \code{0003} ajoute un masque de sondes pour activer individuellement chaque sonde ; les cinq configurations de CT1 sont : \code{plain} (bibliothèque d'origine, aucune sonde), \code{none} (sondes compilées, toutes désactivées), \code{esk} (sonde sur \code{expand_sk_internal} seule), \code{mcalc} (sonde sur \code{m_calculate_M_and_VPV} seule), \code{all} (toutes les sondes).

\section{Analyse statique du code généré (E1)}
Pour chaque construction, l'objet est recompilé par la commande enregistrée, désassemblé (\code{objdump -d -M intel}) et les boucles de travail sont repérées par leur nombre de \code{vpshufb} (36 ymm pour l'officiel de MAYO-5 ; 24 ymm + 6 xmm pour A2). Les comptes sont normalisés par le facteur de déroulage de boucle (ici $u=1$). Le noyau est identifié par le \emph{pas d'adresse sur $P_1$} : \code{add reg,0x47} (entrées de 71 octets, pas d'une entrée) pour \code{P1_times_Vt}, \code{imul ...,0x47} pour \code{P1t_times_Vt} ; les marqueurs de débogage de GCC ne sont pas fiables (les corps identiques sont fusionnés et étiquetés \code{P1_times_O}). Toute attribution d'une différence de code à une cause exige d'avoir écarté les autres modifications du code (section~\ref{sec:e1b}).

\section{Protocole E1b (analyse statique de l'origine des écritures)}\label{sec:e1b}
E1b recompile \code{src/mayo.c} avec \emph{la commande exacte} de E1 plus une seule option supplémentaire, sans aucune mesure de temps. Constructions : \code{P_none} (sans sondes), \code{P_all} (toutes sondes), \code{TAB} (construction des mesures S1/S2). Variantes : \code{base} (contrôle : doit redonner exactement le code de E1) ; \code{dump} (sorties \code{-fdump-tree-optimized}, \code{-fdump-tree-sra-details}, \code{-fdump-tree-esra-details}, \code{-fdump-rtl-expand}, \code{-fdump-rtl-ira}, \code{-fopt-info-inline-optimized} ; contrôle : code identique) ; \code{nosra} (\code{-fno-tree-sra}) ; \code{noipasra} (\code{-fno-ipa-sra}) ; \code{nodse} ; \code{nopre} ; \code{nofre}.

Pour chaque variante, le rapport donne : les boucles de travail avec nombre d'instructions et d'écritures de 32 octets ; le nombre d'instructions de \code{sign_signature_internal} \emph{hors boucles de travail} ; la distance entre histogrammes de mnémoniques hors boucles par rapport à \code{base} ; le nombre de lignes GIMPLE mentionnant \code{temp}. \textbf{Règle d'interprétation :} la disparition des écritures dans une variante n'est jamais lue seule ; elle n'est attribuée à une passe que si la distance hors boucles est faible, c'est-à-dire si rien d'autre n'a changé.

\section{Traçabilité et reproduction}
Les données brutes (S1, S2, PH1, CT1) sont déposées avec leurs empreintes SHA-256 (annexe~\ref{app:inventaire}). Le script de reproduction réextrait les archives, vérifie les empreintes, relance les neuf analyses et compare les sorties : \textbf{9 analyses sur 9 identiques} lors de la dernière exécution. La reproduction des \emph{analyses} est exacte ; celle des \emph{mesures} donnera des valeurs voisines mais non identiques (ticks TSC, fréquence, bruit).

\section{Menaces à la validité identifiées a priori}
\begin{itemize}
\item \textbf{Une machine, un compilateur sur machine physique}, trois démarrages ; S2 sur 5 processus.
\item \textbf{Ticks TSC, non cycles cœur} : la fréquence d'exécution réelle est inconnue ; aucune conclusion sur le débit des unités d'exécution n'est possible.
\item \textbf{Effet d'ordre} documenté (section~\ref{sec:ordre}).
\item \textbf{Les sondes modifient le code} : mesures instrumentées $\ne$ mesures sans sonde.
\item \textbf{Pas de compteurs matériels} disponibles dans l'environnement d'analyse.
\item \textbf{Première exécution de la campagne s3 de S1} : CSV bruts non archivés (seule la seconde exécution l'est).
\item \textbf{Comparaison VM/physique} : le gain de MAYO-5 sur \code{mayo_sign} était de 7,1\,\% dans la VM (GCC 13.3.0) ; l'écart avec la machine physique mêle machine et compilateur et n'est pas expliqué.
\end{itemize}

\chapter{Résultats}
\label{ch:resultats}

Sauf mention contraire, les valeurs sont des rapports A2 SAFE / officiel (médiane des processus), en ticks TSC, sur la machine physique du chapitre~\ref{ch:methodo}. Les statuts \F{}, \I{}, \Hy{} suivent la convention du chapitre~\ref{ch:intro}.

\section{Correction}
\F{} Les KAT officiels de MAYO-1, 2, 3 et 5 réussissent dans toutes les sessions ; 0 signature différente sur les 30\,010 paires de MAYO-5 et les 15\,005 de MAYO-1 et de MAYO-2 ; nombre de tentatives identique dans les deux bras. Les propriétés de sûreté sont au chapitre~\ref{ch:surete}.

\section{Gains sur les noyaux et sur la signature (S1, S2)}
\begin{table}[ht]\centering\small
\caption{Rapport A2/officiel, médiane des processus [IC95 bootstrap entre processus]. S1 : 20 processus ; S2 : 5 processus. Gain = 1 $-$ rapport.}\label{tab:gains}
\begin{tabular}{@{}llccc@{}}\toprule
Paramètre & Niveau & S1 & S2 & Gain S1 / S2\\\midrule
MAYO-1 & \code{P1_times_Vt} & 0,7803 [0,7793 ; 0,7817] & 0,7814 [0,7794 ; 0,7832] & 22,0 / 21,9\,\%\\
 & \code{P1t_times_Vt} & 0,7099 [0,7060 ; 0,7124] & 0,7080 [0,7024 ; 0,7126] & 29,0 / 29,2\,\%\\
 & \code{m_calculate} & 0,8110 [0,8096 ; 0,8118] & 0,8102 [0,8080 ; 0,8123] & 18,9 / 19,0\,\%\\
 & \textbf{\code{mayo_sign}} & \textbf{0,8975} [0,8970 ; 0,8979] & \textbf{0,8974} [0,8968 ; 0,8990] & \textbf{10,3 / 10,3\,\%}\\\midrule
MAYO-5 & \code{P1_times_Vt} & 0,8067 [0,8060 ; 0,8071] & 0,8049 [0,8043 ; 0,8076] & 19,3 / 19,5\,\%\\
 & \code{P1t_times_Vt} & 0,8798 [0,8783 ; 0,8818] & 0,8792 [0,8770 ; 0,8858] & 12,0 / 12,1\,\%\\
 & \code{m_calculate} & 0,8851 [0,8845 ; 0,8863] & 0,8869 [0,8847 ; 0,8874] & 11,5 / 11,3\,\%\\
 & \textbf{\code{mayo_sign}} & \textbf{0,9873} [0,9871 ; 0,9876] & \textbf{0,9871} [0,9863 ; 0,9876] & \textbf{1,3 / 1,3\,\%}\\\bottomrule
\end{tabular}
\end{table}

\begin{figure}[ht]\centering
\begin{tikzpicture}
\begin{axis}[ybar,width=0.92\linewidth,height=6.2cm,bar width=9pt,ymin=0.65,ymax=1.08,
 symbolic x coords={P1,P1t,mcalc,sign},xtick=data,xticklabels={\texttt{P1\_times\_Vt},\texttt{P1t\_times\_Vt},\texttt{m\_calculate},\texttt{mayo\_sign}},x tick label style={font=\scriptsize},
 ylabel={rapport A2 / officiel},legend style={at={(0.5,1.14)},anchor=south,legend columns=2},
 ymajorgrids,nodes near coords,nodes near coords style={font=\tiny,rotate=90,anchor=west},every node near coord/.append style={/pgf/number format/fixed,/pgf/number format/precision=3}]
\addplot[fill=steel!70] coordinates {(P1,0.7803)(P1t,0.7099)(mcalc,0.8110)(sign,0.8975)};
\addplot[fill=navy!70] coordinates {(P1,0.8067)(P1t,0.8798)(mcalc,0.8851)(sign,0.9873)};
\legend{MAYO-1,MAYO-5}
\end{axis}
\end{tikzpicture}
\caption{Rapports A2/officiel mesurés (S1). Le gain se dilue de la couche noyau à la signature complète, fortement pour MAYO-5.}\label{fig:gains}
\end{figure}

\F{} Témoins nuls (aucun noyau A2) : MAYO-2 \code{mayo_sign} 1,0000 / 1,0000, \code{m_calculate} 1,0001 / 0,9996 ; MAYO-3 \code{mayo_sign} 1,0000 / 1,0001. Les différences S2$-$S1 sur les huit lignes de gains sont au plus de 0,0019 en valeur absolue, IC95 contenant 0 (par exemple \code{mayo_sign} : $-0{,}0001$ [$-0{,}0009$ ; $+0{,}0016$] pour MAYO-1, $-0{,}0003$ [$-0{,}0011$ ; $+0{,}0003$] pour MAYO-5).

\paragraph{Variabilité à quatre niveaux.} Écart-type robuste entre paires (dans un processus) : 0,0015 (S1) et 0,0013 (S2) pour MAYO-1, 0,0114 et 0,0039 pour MAYO-5. Entre processus (dans une session) : 0,0009 (20 proc.) et 0,0010 (5 proc.) pour MAYO-1, 0,0008 et 0,0005 pour MAYO-5. Entre campagnes (S1) : 0,0001 et 0,0002. La dispersion entre processus, et non entre paires, fixe donc l'incertitude.

\paragraph{Gain de MAYO-5 sur la signature.} \F{} Il est de 1,3\,\% dans S1, S2, PH1/CT1 et dans la référence, proche du critère de $\approx1$\,\% fixé a priori. C'est le seul chiffre de gain de signature de MAYO-5 à retenir (section~\ref{sec:anomalie}).

\section{Effet d'ordre}\label{sec:ordre}
\F{} Les ordres sont équilibrés : A2 en premier dans 49,97\,\% des paires de \code{P1t_times_Vt} de MAYO-1 en S1 ($29\,993/60\,020$), 50,5\,\% en S2 ; $P(\text{A2 premier}\mid\text{précédent A2})=0{,}5003$.

\begin{table}[ht]\centering\small
\caption{Rapport A2/officiel selon l'ordre d'exécution, et écart par processus [IC95].}\label{tab:ordre}
\begin{tabular}{@{}llcccc@{}}\toprule
Niveau & Session & A2 en 1er & officiel en 1er & écart [IC95] & proc. $>0$\\\midrule
MAYO-1 \code{P1t_times_Vt} & S1 & 0,7142 & 0,6985 & +0,0151 [+0,0101 ; +0,0206] & 20/20\\
 & S2 & 0,7095 & 0,7009 & +0,0043 [$-$0,0070 ; +0,0161] & 4/5\\
MAYO-1 \code{mayo_sign} & S1 & 0,8975 & 0,8975 & $-$0,0001 [$-$0,0002 ; 0,0000] & 8/20\\
MAYO-5 \code{mayo_sign} & S1 & 0,9870 & 0,9876 & $-$0,0006 [$-$0,0008 ; $-$0,0005] & 0/20\\
 & S2 & 0,9871 & 0,9871 & 0,0000 [$-$0,0002 ; 0,0000] & 2/5\\
MAYO-5 \code{P1_times_Vt} & S1 / S2 & 0,8071 / 0,8053 & 0,8059 / 0,8043 & +0,0010 / +0,0010 & 20/20 ; 5/5\\
MAYO-2 \code{m_calculate} (témoin) & S1 / S2 & 1,0005 / 1,0004 & 0,9997 / 0,9994 & +0,0010 / +0,0006 & 20/20 ; 4/5\\\bottomrule
\end{tabular}
\end{table}

\F{} Il existe un effet générique du protocole d'environ 0,1\,\% (visible sur MAYO-2 \code{m_calculate}, où les deux bras sont le même code). L'effet de MAYO-1 \code{P1t_times_Vt} est environ quinze fois plus grand et propre à ce noyau. Il n'affecte pas \code{mayo_sign} de MAYO-1 ($\le0{,}0001$) ; pour MAYO-5 \code{mayo_sign} il vaut au plus $0{,}0006$, soit moins de 5\,\% du gain de $0{,}0127$. Le bras exécuté en second est plus rapide (MAYO-1 \code{P1t_times_Vt}, S1 : officiel $-0{,}9$\,\%, A2 $-1{,}6$\,\%), et l'effet est plus fort sur A2 d'où l'asymétrie du rapport. \Hy{} La cause n'est pas démontrée et aucune expérience dédiée n'a été faite. Pour ce noyau on cite la moyenne des deux ordres (0,706) en plus de la médiane poolée (0,710).

\section{L'anomalie de MAYO-5 : décomposition instrumentée et contrôle}\label{sec:anomalie}
\subsection{Toutes sondes actives (PH1)}
\F{} MAYO-5 : une tentative de signature par signature dans les deux bras ; chaque phase est appelée une fois par signature, dont \code{P1_times_Vt} et \code{P1t_times_Vt} une fois par tentative. Temps total 1\,374\,501 (officiel) et 1\,274\,140 (A2) : rapport 0,9270 [0,9253 ; 0,9283], soit 7,3\,\%. La différence moyenne de $-100\,967$ ticks se répartit en $-100\,529$ pour \code{m_calculate} (99,6\,\%), dont \code{P1_times_Vt} $-54\,584$ et \code{P1t_times_Vt} $-46\,166$ ; les autres phases (dont \code{expand_sk_internal}, \code{sample_solution}, \code{compute_A}) sont identiques à $\pm0{,}1$\,\%. MAYO-1 : rapport 0,8983 ; MAYO-2 (témoin nul) : 0,9999.

\subsection{Le gain avec toutes les sondes n'est pas le gain sans sonde (CT1)}
\begin{table}[ht]\centering\small
\caption{Cinq configurations du même programme (10 processus chacune, graines 41--50, un démarrage). Ticks par signature.}\label{tab:ct1}
\begin{tabular}{@{}lrrcc@{}}\toprule
Configuration & Officiel & A2 & Rapport [min--max] & Gain\\\midrule
\code{plain} (aucune sonde) & 1\,504\,741 & 1\,484\,504 & 0,9871 [0,9862--0,9879] & 1,29\,\%\\
\code{none} (sondes compilées, désactivées) & 1\,504\,802 & 1\,484\,611 & 0,9868 [0,9860--0,9880] & 1,32\,\%\\
\code{esk} (\code{expand_sk_internal}) & 1\,504\,167 & 1\,484\,146 & 0,9871 [0,9861--0,9879] & 1,29\,\%\\
\code{mcalc} (\code{m_calculate_M_and_VPV}) & 1\,502\,948 & 1\,483\,529 & 0,9877 [0,9865--0,9885] & 1,23\,\%\\
\code{all} (toutes les sondes) & 1\,372\,740 & 1\,273\,970 & 0,9276 [0,9261--0,9291] & 7,24\,\%\\\bottomrule
\end{tabular}
\end{table}

\begin{figure}[ht]\centering
\begin{tikzpicture}
\begin{axis}[ybar,width=0.85\linewidth,height=5.2cm,bar width=16pt,ymin=0.90,ymax=1.0,
 symbolic x coords={plain,none,esk,mcalc,all},xtick=data,ylabel={rapport A2 / officiel},ymajorgrids,
 nodes near coords,every node near coord/.append style={font=\scriptsize,/pgf/number format/fixed,/pgf/number format/precision=4}]
\addplot[fill=steel!70] coordinates {(plain,0.9871)(none,0.9868)(esk,0.9871)(mcalc,0.9877)(all,0.9276)};
\end{axis}
\end{tikzpicture}
\caption{MAYO-5, rapport \texttt{mayo\_sign} A2/officiel selon la configuration de sondes (CT1). Seule la configuration « toutes les sondes » change le résultat.}\label{fig:ct1}
\end{figure}

\subsection{Mesure directe, sans soustraction de médianes (config. \code{mcalc})}
\begin{table}[ht]\centering\small
\caption{Décomposition directe de la signature MAYO-5 (ticks par signature).}\label{tab:direct}
\begin{tabular}{@{}lrrr@{}}\toprule
 & Officiel & A2 & Différence\\\midrule
\code{m_calculate} dans la signature & 931\,453 & 912\,476 & $-18\,406$ [$-19\,396$ ; $-17\,558$]\\
hors \code{m_calculate} (total $-$ sonde) & 566\,917 & 567\,323 & $+406$ (0,07\,\%)\\
\code{expand_sk_internal} (config. \code{esk}) & 322\,655 & 322\,661 & $+6$\\\midrule
\code{m_calculate} isolé (banc, tampons réutilisés) & 801\,324 & 707\,314 & $\approx-94\,000$ (11,7\,\%)\\
\code{m_calculate}, toutes sondes, dans la signature & 801\,073 & 699\,775 & $\approx-101\,000$\\\bottomrule
\end{tabular}
\end{table}

\F{} Le reste de la signature est identique dans les deux versions. La différence est entièrement dans \code{m_calculate} \emph{en contexte de signature sans sondes internes} : il coûte 931\,000 (officiel) et 912\,000 (A2) ticks, contre environ 801\,000 et 707\,000 isolé ou sondé. Le surcoût de contexte est d'environ 130\,000 ticks (officiel) et 205\,000 (A2), et la différence de 75\,000 ticks entre les deux absorbe l'essentiel du gain isolé. Il disparaît avec toutes les sondes (donc des \code{lfence} entre les sous-étapes de \code{m_calculate}) mais pas avec une sonde unique autour de \code{m_calculate}.

\paragraph{Correction d'une interprétation antérieure.} Une première lecture (« le reste de la signature ralentit avec A2 ») résultait d'une soustraction de médianes issues de niveaux de mesure différents ; elle est \emph{retirée}.

\section{Analyse statique du code généré (E1, GCC~16.2.0)}
\subsection{Comptage de base}
\F{} Par entrée de $P_1$, normalisé ($u=1$) : \code{vpshufb} officiel 20 (MAYO-1) et 36 (MAYO-5), A2 15 (10 ymm + 5 xmm) et 30 (24 ymm + 6 xmm), conformément au modèle du chapitre~\ref{ch:a2}. Dans le banc, instructions de la boucle interne : MAYO-5 \code{P1_times_Vt} 124 $\to$ 106, \code{P1t_times_Vt} 126 $\to$ 110 ; MAYO-1 \code{P1_times_Vt} 68 $\to$ 51, \code{P1t_times_Vt} 71 $\to$ 56. Rechargements de pile : 30 $\to$ 24 (MAYO-5 \code{P1_times_Vt}) et 11 $\to$ 5 (MAYO-1). Cette réduction est une explication appuyée sur l'assembleur, pas une preuve que ces instructions expliquent à elles seules le gain.

\subsection{Écritures supplémentaires en contexte de signature}
\begin{table}[ht]\centering\small
\caption{Instructions de la boucle interne par entrée de $P_1$ (MAYO-5) selon la construction.}\label{tab:e1}
\begin{tabular}{@{}lcccc@{}}\toprule
Construction & \code{P1} off. & \code{P1} A2 & \code{P1t} off. & \code{P1t} A2\\\midrule
banc isolé & 124 & 106 & 126 & 110\\
contexte, toutes sondes & 123 & 105 & 126 & 109\\
contexte, sans toutes les sondes (\code{ctx_ab}, \code{esk}, \code{mcalc}, \code{none}) & 123 & 135 & 162 & 139\\\bottomrule
\end{tabular}
\end{table}

\F{} La différence entre « sans toutes les sondes » et « toutes les sondes » est \emph{exactement} une écriture de 32 octets (\code{vmovdqa [rbp-...], ymm}) de plus par accumulateur et par entrée : +36 pour \code{P1t_times_Vt} officiel (36 accumulateurs), +30 pour chaque noyau A2 (30 accumulateurs). Le reste de la boucle est identique instruction pour instruction. La boucle officielle \code{P1_times_Vt} n'a pas de doublon. Les écritures vont vers une seconde série d'emplacements consécutifs (pas de \code{0x20}), distants de \code{0xbfc0} octets de la première, dont l'adresse est prise ailleurs dans la fonction (\code{lea}). Deux boucles de 160 instructions (\code{Ut_times_O}, \code{m_lambda_times_O}) portent aussi un doublon (123 $\to$ 160) dans \emph{les deux} bras. Le cadre de pile dépasse 220 Ko, ce qui interdit de distinguer par le seul décalage un débordement de registres d'un tableau local. Avec GCC~13.3 (VM), aucun doublon n'apparaît. Pour MAYO-1 (GCC~16.2.0), le doublon touche seulement \code{P1t_times_Vt} (officiel 71 $\to$ 91 ; A2 56 $\to$ 71).

\F{} L'identification de la boucle officielle \code{P1_times_Vt} dans le contexte repose sur le pas d'adresse \code{add reg,0x47} et a été contrôlée dans six constructions ; elle donne 123 instructions et 29 écritures dans les contextes officiel, sans sondes et avec toutes les sondes (adresses distinctes), sans doublon. Le tableau~\ref{tab:e1} résume ces comptes.

\section{Lien quantitatif entre instructions et temps : modèle à un paramètre}
\F{} Gain statique (P1+P1t) : banc 13,6\,\%, contexte toutes sondes 14,1\,\%, contexte sans sondes \textbf{3,9\,\%} (250 $\to$ 216 ; 249 $\to$ 214 ; 285 $\to$ 274 instructions). Gain mesuré de \code{m_calculate} : 11,5--11,7\,\% isolé, 12,4\,\% toutes sondes, \textbf{2,1\,\%} sans sondes (936\,031 $\to$ 916\,206). Dans les trois constructions, l'avantage statique et l'avantage mesuré évoluent dans le même sens.

\paragraph{Modèle.} Le temps d'une boucle est supposé proportionnel à son nombre d'instructions par entrée avec un coefficient $c$ commun aux deux bras ($N_{\mathrm{ent}}=10\,153$). Étalonné sur « toutes sondes » : $c=635\,681/(249\cdot10\,153)=0{,}251$ tick/instruction. \textbf{Prédiction} sans sondes : gain de P1+P1t de $28\,082$ ticks ; \textbf{mesuré} : $19\,825$. Le modèle \emph{surestime} le gain d'environ 40\,\%.

\paragraph{Coût par écriture supplémentaire.} À partir de la différence « sans sondes » $-$ « toutes sondes » : officiel $130\,895/365\,508=0{,}358$ ; A2 $210\,808/609\,180=0{,}346$ tick par écriture.

\begin{encadre}
\textbf{Portée de ce rapprochement.} La proximité de 0,358 et de 0,346 est un simple test de cohérence (un paramètre, deux observations) ; \emph{elle ne constitue pas une validation du modèle}. Le modèle surestimant le gain, il n'explique pas quantitativement l'écart. La corrélation entre écritures supplémentaires et perte de gain n'est pas une preuve de causalité.
\end{encadre}

\section{Analyse statique E1b : état}
Le protocole (section~\ref{sec:e1b}) est implémenté (\code{e1b.py}) et validé mécaniquement sur une machine de test avec GCC~13 (contrôle de reproduction du code de E1 : identique ; distance des histogrammes hors boucles : 0 pour la variante \code{dump}, 2 pour \code{nosra}). \textbf{Il n'a pas encore été exécuté avec GCC~16.2.0 sur la machine de mesure ; aucun résultat d'E1b n'est donc présenté.} La section~\ref{sec:lecture-e1b} fixe à l'avance la lecture des issues possibles.

\subsection{Lecture prévue des issues (fixée avant exécution)}\label{sec:lecture-e1b}
\begin{table}[ht]\centering\small
\caption{Issues possibles d'E1b et conclusions autorisées.}\label{tab:e1b-issues}
\begin{tabularx}{\linewidth}{@{}XXX@{}}\toprule
Observation (GCC~16.2.0) & Conclusion autorisée & Conclusion \emph{non} autorisée\\\midrule
Doublon présent dans \code{base} et \code{dump}, absent avec \code{-fno-tree-sra} \emph{et} distance hors boucles faible & la passe SRA conditionne le doublon dans cette compilation (à confirmer par les dumps) & que \code{temp[]} en soit la cause tant que les dumps GIMPLE ne le montrent pas\\
Doublon absent avec \code{-fno-tree-sra} mais distance hors boucles élevée & la désactivation change aussi d'autres parties du code ; l'attribution à SRA est \emph{non établie} & toute attribution à SRA\\
Doublon présent dans toutes les variantes & les passes testées ne le suppriment pas ; l'hypothèse « SRA » n'est pas confirmée & que le doublon soit sans cause compilateur\\
\code{base} $\ne$ code de E1 & E1 non reproduit : arrêter l'interprétation & --\\\bottomrule
\end{tabularx}
\end{table}

\section{Comparaison avec la VM}
\F{} Dans la VM (Xeon virtualisé, GCC~13.3.0), le gain de \code{mayo_sign} pour MAYO-5 était de 7,1\,\%, proche du 7,3\,\% mesuré avec toutes les sondes sur la machine physique. \Hy{} L'écart avec 1,3\,\% sans sonde mêle machine et compilateur ; il n'est pas expliqué. Ce rapprochement est mentionné sans être exploité.

\chapter{Discussion critique}
\label{ch:discussion}

Ce chapitre sépare ce qui est établi, ce qui est inféré et ce qui reste hypothèse, puis discute les limites et les pistes écartées.

\section{Bilan épistémique}
\begin{table}[ht]\centering\small
\caption{Établi, inféré, hypothétique.}\label{tab:bilan}
\begin{tabularx}{\linewidth}{@{}lXl@{}}\toprule
Statut & Énoncé & Appui\\\midrule
\F{} & Nombre de \code{vpshufb} par entrée : 20$\to$15 (MAYO-1), 36$\to$30 (MAYO-5) & assembleur, preuve (prop.~\ref{prop:real})\\
\F{} & Gains de noyaux 12--29\,\% ; \code{mayo_sign} 10,3\,\% (MAYO-1), 1,3\,\% (MAYO-5) sur cette machine et ce compilateur & S1, S2, PH1, CT1\\
\F{} & Le reste de la signature est identique entre bras ; la différence est dans \code{m_calculate} en contexte & mesure directe (\code{mcalc})\\
\F{} & Avec toutes les sondes, le gain de MAYO-5 est 7,3\,\% ; sans, 1,3\,\% & CT1\\
\F{} & GCC~16.2.0 émet une écriture de 32 octets supplémentaire par accumulateur en contexte sans toutes les sondes, asymétrique entre bras & désassemblage E1\\
\I{} & L'asymétrie de ces écritures (+60 pour A2, +36 pour l'officiel sur P1+P1t) explique l'essentiel de la perte de gain de MAYO-5 en contexte & concordance statique/mesure\\
\Hy{} & Le coût passe par le débit des écritures plutôt que par les dépendances écriture--lecture ou conflits d'adresses & non testé\\
\Hy{} & Le tableau \code{temp[]} indexé par la variable de boucle est conservé en mémoire en plus des emplacements scalaires & E1b à exécuter\\
\Hy{} & Les sondes modifient la structure du code et donc la décision de GCC & cohérent avec E1, non démontré\\
\bottomrule
\end{tabularx}
\end{table}

\section{Pourquoi le gain se dilue}
Le gain d'un noyau ne se reporte sur la signature qu'au prorata de son poids. Pour MAYO-5 sans sondes, \code{m_calculate} pèse environ $931\,453/1\,504\,741\approx62$\,\% de la signature et son gain en contexte est de $18\,406$ ticks, soit $1{,}2$\,\% de la signature : c'est, à la précision des mesures, le gain de 1,3\,\% observé. Pour MAYO-1, le rapport entre gain de signature (10,3\,\%) et gain de \code{m_calculate} (18,9\,\%) donne un poids implicite d'environ 54\,\% \I{}. La question n'est donc pas la fraction de la signature occupée par \code{m_calculate}, mais pourquoi son gain tombe de 11,7\,\% (isolé) à 2,0\,\% (en contexte sans sondes) pour MAYO-5, alors qu'il reste de 18,9\,\% pour MAYO-1 \F{}.

\section{Ce que la concordance statique/dynamique prouve et ne prouve pas}
Le désassemblage donne un fait solide : le nombre d'instructions de la boucle diffère selon la construction, de façon asymétrique. Les temps varient dans le même sens. Trois précautions s'imposent.
\begin{enumerate}
\item \textbf{Un modèle à un paramètre qui surestime de 40\,\%} le gain ne peut pas servir de validation ; son accord sur le coût par écriture entre les deux bras (0,358 contre 0,346) est un contrôle de cohérence, non une preuve.
\item \textbf{Corrélation, non causalité} : les constructions qui diffèrent par le doublon diffèrent aussi par la présence de \code{lfence} entre sous-étapes, par la taille de la fonction et par l'allocation de la pile. Seule une expérience qui supprime le doublon \emph{sans autre changement de code} permettrait de conclure (cf.\ la conclusion générale).
\item \textbf{Ticks et non cycles} : le « coût par écriture » de 0,35 tick équivaut, à une fréquence soutenue non mesurée de 3,5--4,5 GHz pour un TSC de 2,4 GHz, à un ordre de grandeur de 0,5--0,65 cycle ; la compatibilité avec un débit d'une écriture par cycle est une remarque \Hy{}, pas un résultat.
\end{enumerate}
Plusieurs raisons possibles d'une surestimation du modèle sont envisageables (toutes \Hy{}) : le coût des instructions n'est pas homogène (un \code{vpshufb} et une écriture n'ont pas le même coût) ; des boucles hors P1/P1t portent aussi un doublon (\code{Ut_times_O}, \code{m_lambda_times_O}) dont les nombres de tours ne sont pas connus ici ; les chargements sur la même ligne de cache peuvent se recouvrir.

\section{Pourquoi MAYO-1 et MAYO-5 se comportent différemment}
\I{} Pour MAYO-1, le doublon ne touche que \code{P1t_times_Vt} dans les deux bras à un rapport voisin (71/91 $=0{,}78$ contre 56/71 $=0{,}79$ dans le banc) : l'avantage statique est conservé (0,770 banc, 0,774 contexte), en accord avec un gain de signature qui ne s'effondre pas. Pour MAYO-5, le doublon touche les deux noyaux d'A2 mais un seul de l'officiel. Cette lecture est compatible avec les faits mais elle reste une inférence, pour la raison donnée ci-dessus.

\section{Pistes explorées puis écartées}
Les pistes suivantes ont été explorées dans la VM avant la machine physique ; elles sont consignées comme hypothèses éliminées (avec leurs limites).

\paragraph{Fusion de \code{P1_times_Vt} et \code{P1t_times_Vt} en un seul parcours.} Idée : lire $P_1$ une fois au lieu de deux. Le nombre de lectures d'entrées est bien divisé par deux (MAYO-1 : 6\,480 $\to$ 3\,240 ; MAYO-5 : 20\,306 $\to$ 10\,153) mais, avec accumulateurs simultanés (option 1), le temps de \code{mayo_sign} augmente de 4,9 à 10,4\,\% (MAYO-1) et de 5,3 à 7,0\,\% (MAYO-5) ; avec fermeture immédiate (option 2), de 67 à 78\,\% (MAYO-1, MAYO-3, MAYO-5). L'option 1 exige $V_{\max}$ accumulateurs, soit +49 Kio (MAYO-1) à +159 Kio (MAYO-5) de pile. Correction vérifiée (KAT et 5\,000 essais indépendants par paramètre sans écart). Interprétation : le travail SIMD est inchangé et la pression mémoire des accumulateurs dispersés l'emporte sur la réduction des lectures. Piste close ; mesures dans la VM uniquement.

\paragraph{Transposée de $P_1$.} Un diagnostic sur la représentation et le motif d'accès à $P_1$ conclut que la disposition joue un rôle important dans les 29--31\,\% de gain théorique maximal observés sur MAYO-3 et MAYO-5, sans attribution quantitative au cache (pas de compteurs matériels), et que le coût de construction l'emporte sur le gain. Piste close.

\paragraph{Traitement par lots dans EF.} Piste fermée ; patch, banc et données archivés. Une anomalie de dérive entre processus sur MAYO-2 a été documentée : l'amplitude d'un rapport à un seul niveau dépend du régime du processus, ce qui impose de lire les rapports avec la dérive de la phase correspondante.

\section{Limites de la validation}
Une seule machine ; un seul compilateur sur la machine physique ; trois démarrages ; S2 sur cinq processus ; gain de MAYO-5 proche du critère de 1\,\% ; sondes complètes qui modifient le comportement ; pas de cycles cœur ; pas de compteurs matériels ; Clang non examiné ; aucune preuve formelle de temps constant ; les propriétés de bourrage reposent sur une énumération des consommateurs de \code{Pv}. Par ailleurs l'article n'a été lu intégralement que dans les parties pertinentes (§5.1 et code AVX2).

\section{Sur l'intégration dans MAYO-C}
Aucune décision d'intégration n'est prise ici. Les éléments à peser sont : le gain de 10,3\,\% pour MAYO-1 (robuste sur deux démarrages) ; le gain de 1,3\,\% pour MAYO-5 (proche du critère, dépendant du code généré) ; les préconditions (voies égales, $m\in\{80,142\}$) ; l'état de la preuve de temps constant (indicateurs seulement) ; la dépendance au compilateur. Un jugement d'intégration plus solide exigerait au moins une seconde machine ou un second compilateur.

\chapter*{Conclusion générale et perspectives}
\addcontentsline{toc}{chapter}{Conclusion générale et perspectives}
\markboth{Conclusion générale et perspectives}{Conclusion générale et perspectives}

\section*{Contributions effectives}
Ce mémoire a étudié, dans MAYO-C, la façon dont les noyaux AVX2 \code{P1_times_Vt} et \code{P1t_times_Vt} traitent la fin des $m$-vecteurs. Les contributions établies sont les suivantes.
\begin{enumerate}
\item Un \textbf{comptage rigoureux} dans un modèle précis : le minimum est $\lceil m/32\rceil$ instructions \vpshufb{}, le schéma officiel en utilise $2\lceil m/64\rceil$, et un gain n'est possible que si $0<m\bmod64\le32$. MAYO-1 et MAYO-5 en profitent, pas MAYO-2 ni MAYO-3.
\item Une \textbf{mise en œuvre} (A2) pour une fin de vecteur d'au plus 16 éléments, avec une variante SAFE qui ne lit rien au-delà de l'entrée de $P_1$. Les propriétés de lecture, d'écriture, de tables et de bourrage sont justifiées par le code et vérifiées par des tests.
\item Des \textbf{gains mesurés} sur une machine physique, avec une unité statistique précisée : 12 à 29\,\% sur les noyaux, 10,3\,\% sur \code{mayo_sign} pour MAYO-1 et 1,3\,\% pour MAYO-5, retrouvés lors de deux démarrages de la machine.
\item La \textbf{localisation} de la perte de gain de MAYO-5 dans \code{m_calculate}, quand on mesure la signature complète. On observe aussi, dans le code produit par GCC~16.2.0, des écritures mémoire en plus et inégales entre les deux versions, qui correspondent à cette perte. L'analyse E1b montre que ce sont des stockages de tableaux d'accumulateurs refaits à chaque tour de boucle, et écarte la passe SRA comme cause.
\item Un \textbf{dépôt reproductible} (scripts, données brutes, empreintes) avec lequel les neuf analyses donnent exactement les mêmes résultats.
\end{enumerate}

\section*{Limites}
Le gain sur la signature de MAYO-5 est faible et dépend du choix du compilateur. Le lien de cause à effet entre les écritures en plus et la perte de gain n'est pas démontré, et le modèle chiffré exagère l'effet d'environ 40\,\%. Les mesures viennent d'une seule machine et d'un seul compilateur, en ticks TSC. Le temps d'exécution indépendant du secret n'est appuyé que par des indices, pas par une preuve. La variante pour une fin de vecteur de 17 à 32 éléments n'est pas réalisée. Nous ne prétendons à aucune nouveauté absolue.

\section*{Questions ouvertes et perspectives}\label{sec:persp}
\begin{description}
\item[Q1 -- Le choix du compilateur.] Comprendre pourquoi GCC~16.2.0 laisse ces écritures dans la boucle quand toutes les sondes ne sont pas actives, et les sort de la boucle quand elles le sont. L'analyse E1b a écarté SRA, DSE, PRE et IPA-SRA comme causes uniques (en les désactivant une à une, le doublon reste) et ne permet pas d'accuser FRE (le code change aussi ailleurs). La suite naturelle, toujours sans exécuter le programme, est de produire les fichiers d'analyse des passes qui déplacent les stockages et de PRE (\code{-fdump-tree-lim2-details}, \code{-fdump-tree-pre-details}). Ce n'a pas encore été fait : cela attend l'accord de l'encadrant.
\item[Q2 -- La cause.] Si l'on trouve un moyen de supprimer le doublon sans toucher à autre chose dans le code, comparer alors les temps avec et sans doublon, sur le même code. Les expériences E2 à E5 prévues au départ restent en attente de la décision de l'encadrant.
\item[Q3 -- La solidité des résultats.] Refaire l'analyse avec Clang et d'autres versions de GCC, et mesurer sur une deuxième machine, si possible avec les compteurs du processeur et en cycles d'horloge (hypothèse de performance H-P).
\item[Q4 -- L'extension.] Réaliser et valider la variante pour $16<r\le32$ ; examiner d'autres paramètres que ceux étudiés ici.
\item[Q5 -- La garantie.] Vérifier, de façon formelle ou avec un outil, que le programme compilé s'exécute en temps indépendant du secret.
\item[Q6 -- L'intégration.] Décider, en fonction des réponses précédentes et des critères de l'encadrant.
\end{description}


\appendix
\chapter{Inventaire des archives et reproduction}\label{app:inventaire}
État au 22 octobre 2026 (dernier inventaire vérifié) ; les ajouts postérieurs (E1b, identification des noyaux) sont dans le dépôt \code{mayo-a2-thesis}. Toutes les empreintes sont des SHA-256.

\section{Versions définitives à citer}
\begingroup\small\setlength{\tabcolsep}{3pt}
\begin{longtable}{@{}p{0.30\linewidth}p{0.42\linewidth}p{0.22\linewidth}@{}}
\toprule Fichier & SHA-256 & Contenu\\\midrule\endhead
\code{A2_partie1.tar.gz} & \sha{cbc0cc297ee0a0790b3cb7513d0384c6c323cfe26b6180b93fc7bd1b115477a6} & notes, code, scripts, analyses, brut S2\\
\code{A2_partie2.tar.gz} & \sha{d8dbf1b169f249c720a6e74eb0b9c6e6c9fe44eae6087ef6fe188aa21ab0fcab} & brut PH1 et CT1\\
\code{A2_partie3.tar.gz} & \sha{00d12e885de273e7851b270f8b53b780eb52fd5bbbe969b1dfe73402b687ac33} & brut S1\\
patch \code{0001} & \sha{e6c03362cbf974cf750f340453de874d667b5670a7e60eddafa54b35e29e0cea} & prototype A2\\
patch \code{0002} & \sha{751ec21fbd059e51a4b54f193f3e346a889d4b361cfdf78a3ab77c3e891b6280} & sondes\\
patch \code{0003} & \sha{b590a81d479fd4a32883a60efe4bb5525fb9f45bf92a1acd70f60b265596112d} & contrôle des sondes\\
patch \code{0004} & \sha{9c460d0de2193805f3b0c1716f62bf1390ca543778ccfb0022d995ef254a10c6} & analyse statique E1\\
{\ttfamily\scriptsize MAYO-C-A2\_source\_ e72d407.tar.gz} & \sha{e2252a0b51ab827bd097e5a0df73dcbb7ca13539af82b1e728737c58800f6833} & arbre S1/S2\\
{\ttfamily\scriptsize MAYO-C-A2-phases-v2\_ source.tar.gz} & \sha{ecb0941b086affb7189151af4b58355acb39edd4de8d7089b34ac7328cf19ef3} & arbre PH1/CT1\\
\code{E1_kit.tar.gz} & \sha{efe0b71ca5dc1553b0c0a841dc59096216a214080d988ba32731cbe2fcde18f3} & \code{run_E1.sh}, \code{e1_analyze.py}\\
\bottomrule
\end{longtable}
\endgroup
Archive E1b (sorties de compilation, dumps GCC, désassemblages) : \code{e1b_gcc16.tar.gz}, SHA-256 \sha{49452120b1b787ad2aed0f30cbcece4f7584a6063768757d92214e264dcd3faa}.\par
Archives brutes internes : S1 \code{ebc01dfe}\dots, S2 \code{b1a69769}\dots c723, PH1 \code{e1a88aec}\dots 98cc, CT1 \code{0506c2b5}\dots 92de. Les archives de la phase VM (\code{A2_donnees_et_journaux}, \code{tail_aware_etude}, \code{archive_piste_batch_EF}, etc.) sont de l'\emph{historique} et ne doivent jamais être présentées comme performances matérielles.

\section{Identification exacte du code}
Le code est déterminé par l'arbre git (indépendant de l'horodatage des commits). Base : \code{PQCMayo/MAYO-C} commit \code{74de637a7c4a7f9fafa888183bc05f5cf920cd2d}, puis \code{git am} des patchs dans l'ordre. Arbres attendus : après \code{0001} \code{87d1729c}\dots (S1, S2) ; après \code{0002} \code{c8c37c97}\dots ; après \code{0003} \code{92bcede1}\dots (PH1, CT1) ; après \code{0004} \code{bf6955da}\dots.
\begin{verbatim}
git clone https://github.com/PQCMayo/MAYO-C && cd MAYO-C
git checkout 74de637a7c4a7f9fafa888183bc05f5cf920cd2d
git am 0001-*.patch 0002-*.patch 0003-*.patch 0004-*.patch
git rev-parse HEAD^{tree}   # attendu : bf6955dae657abb087b58e13e314b946e0d283fe
\end{verbatim}

\section{Reproduire les analyses}
Prérequis : Linux, \code{python3} avec \code{numpy}, \code{tar}, \code{sha256sum}. Les trois archives dans un dossier $D$ :
\begin{verbatim}
sh reproduire_analyses.sh D travail_neuf
\end{verbatim}
Le script vérifie les empreintes, relance les neuf analyses et affiche \code{OK} si la sortie est identique (résultat obtenu : 9 sur 9). Les sorties couvrent la synthèse S1/S2 (\code{synthese_sessions.py}), la comparaison (\code{compare_sessions.py}), la décomposition (\code{phases_analyze.py}), le contraste avec/sans sondes (\code{phases_vs_ref.py}) et le contrôle CT1 (\code{controle_analyze.py}).

\section{Refaire les mesures}
Conditions : Linux, AVX2, GCC $\ge13$ (mesures : 16.2.0), cmake, \code{python3}+\code{numpy}, gouverneur \code{performance} (\code{sudo cpupower frequency-set -g performance}), cœur 1 dédié, aucun autre travail pendant l'exécution.
\begin{center}\small\begin{tabular}{@{}ll@{}}\toprule
Étape & Commande (racine de l'arbre correspondant)\\\midrule
S1, S2 & \code{SESSION=<étiquette> sh run_session_independante.sh}\\
PH1 & \code{SESSION=<étiquette> sh test/run_phases.sh}\\
CT1 & \code{SESSION=<étiquette> sh test/run_controle.sh}\\
E1 & \code{SESSION=<étiquette> sh test/run_E1.sh}\\
E1b & \code{python3 -I e1b.py <dossier_E1> <dossier_neuf>}\\\bottomrule
\end{tabular}\end{center}
Chaque \code{machine.txt} enregistre date UTC, noyau, compilateur, microcode, gouverneur, \code{boot_id} et alimentation. La reproduction des \emph{mesures} donnera des valeurs voisines et non identiques.

\chapter{Identification des noyaux dans le code généré (GCC~16.2.0, MAYO-5)}\label{app:ident}
Identification par le pas d'adresse sur $P_1$ (section~\ref{sec:protocole}). « Écritures » : écritures de 32 octets dans le corps de boucle.

\begin{table}[ht]\centering\small
\caption{Instructions / écritures de 32 octets des boucles de travail selon la construction.}
\begin{tabular}{@{}lcccc@{}}\toprule
Construction & \code{P1} off. & \code{P1} A2 & \code{P1t} off. & \code{P1t} A2\\\midrule
\code{ctx_official} & 123 / 29 & -- & 162 / 65 & --\\
\code{ctx_ab}, \code{none}, \code{esk}, \code{mcalc} & 123 / 29 & 135 / 52 & 162 / 65 & 139 / 52\\
\code{ctx_phases} (toutes sondes) & 123 / 29 & 105 / 22 & 126 / 29 & 109 / 22\\\bottomrule
\end{tabular}
\end{table}
Les différences « sans toutes les sondes » moins « toutes sondes » sont : A2 \code{P1} $+30$ (52$-$22), A2 \code{P1t} $+30$, officiel \code{P1t} $+36$ (65$-$29), officiel \code{P1} $0$.


\chapter{Exemple jouet complet de MAYO sur $\mathbb{F}_{16}$}\label{app:jouet}
Cet exemple, calculé à la main et vérifié par un programme, illustre la section~\ref{sec:schema}. Les éléments de $\mathbb{F}_{16}=\mathbb{F}_2[x]/(x^4+x+1)$ sont écrits en hexadécimal (A $=10$, B $=11$, \dots, F $=15$) ; l'addition est le OU exclusif.

\section*{Paramètres}
$n=3$, $v=2$, $o=1$, $m=2$, $k=2$ : on a bien $o<m$ et $k\,o=m$. Le secret est $O=(\mathrm{B},8)^{\mathsf T}$.

\section*{Clé publique (blocs de chaque équation $\ell=1,2$)}
\[P^{(1)}_1=\begin{pmatrix}8&9\\0&6\end{pmatrix},\ P^{(1)}_2=\begin{pmatrix}2\\2\end{pmatrix},\ P^{(1)}_3=(6);\qquad
P^{(2)}_1=\begin{pmatrix}8&\mathrm{A}\\0&8\end{pmatrix},\ P^{(2)}_2=\begin{pmatrix}\mathrm{B}\\\mathrm{C}\end{pmatrix},\ P^{(2)}_3=(\mathrm{E}).\]
Le bloc $P^{(\ell)}_3=\mathrm{Upper}\bigl(-O^{\mathsf T}(P^{(\ell)}_1O-P^{(\ell)}_2)\bigr)$ est calculé avec le secret.

\section*{Signature}
Le haché salé est supposé égal à $t=(\mathrm{F},9)$. Vinaigres tirés : $V=\begin{pmatrix}2&3\\3&\mathrm{A}\end{pmatrix}$ (lignes $v_1,v_2$). On calcule
\[Y^{(1)}=VP^{(1)}_1V^{\mathsf T}=\begin{pmatrix}8&7\\6&4\end{pmatrix},\qquad Y^{(2)}=VP^{(2)}_1V^{\mathsf T}=\begin{pmatrix}1&4\\3&9\end{pmatrix}.\]
Le second membre est $y=t+\bigl(u_{11}+z\,u_{12}+z^2u_{22}\bigr)=(8,7)$, où $z$ est la classe de $x$ modulo $z^2+z+1$ et $u_{ij}$ le $m$-vecteur des coefficients de $Y^{(\cdot)}$. La matrice du système est $A=\begin{pmatrix}2&1\\\mathrm{F}&\mathrm{D}\end{pmatrix}$ ; son déterminant vaut $6\neq0$, donc il y a une solution unique $x=(x_1,x_2)=(8,\mathrm{B})$. La signature est $S=(s_1,s_2)$ avec $s_i=(v_i+Ox_i\,\|\,x_i)$ :
\[s_1=(5,\mathrm{F},8),\qquad s_2=(\mathrm{A},\mathrm{D},\mathrm{B}).\]

\section*{Vérification}
On évalue $P^{(\ell)}(s_1)$, $P^{(\ell)}(s_2)$ et $P'^{(\ell)}(s_1,s_2)$ :
$P(s_1)=(\mathrm{A},5)$, $P(s_2)=(9,\mathrm{B})$, $P'(s_1,s_2)=(\mathrm{D},4)$. Alors
\[P^\star(S)=P(s_1)+z\,P'(s_1,s_2)+z^2P(s_2)=(\mathrm{A},5)+(6,9)+(3,5)=(\mathrm{F},9)=t.\]
La signature est acceptée. Si l'on modifie un seul élément de $s_1$, l'égalité est perdue et la signature est refusée.

\begin{thebibliography}{99}
\addcontentsline{toc}{chapter}{Bibliographie}
\bibitem{shor} P. W. Shor, \emph{Algorithms for quantum computation: discrete logarithms and factoring}, FOCS 1994.
\bibitem{patarin} J. Patarin, \emph{The oil and vinegar signature scheme}, Dagstuhl Workshop on Cryptography, 1997.
\bibitem{kipnisshamir} A. Kipnis, A. Shamir, \emph{Cryptanalysis of the Oil and Vinegar Signature Scheme}, CRYPTO 1998.
\bibitem{uov} A. Kipnis, J. Patarin, L. Goubin, \emph{Unbalanced Oil and Vinegar Signature Schemes}, EUROCRYPT 1999.
\bibitem{beullensrainbow} W. Beullens, \emph{Breaking Rainbow Takes a Weekend on a Laptop}, CRYPTO 2022.
\bibitem{mayo} W. Beullens, \emph{MAYO: Practical Post-Quantum Signatures from Oil-and-Vinegar Maps}, IACR ePrint 2021/1144 (SAC 2021).
\bibitem{nibbling} W. Beullens, F. Campos, S. Celi, B. Hess, M. J. Kannwischer, \emph{Nibbling MAYO: Optimized Implementations for AVX2 and Cortex-M4}, IACR TCHES 2024(2).
\bibitem{nistfips} NIST, \emph{FIPS 203, 204, 205: normes de chiffrement et de signature post-quantiques}, août 2024.
\bibitem{nistadd} NIST, \emph{Call for Additional Digital Signature Schemes for the Post-Quantum Cryptography Standardization Process}, 2022.
\bibitem{kocher} P. Kocher, \emph{Timing Attacks on Implementations of Diffie-Hellman, RSA, DSS, and Other Systems}, CRYPTO 1996.
\bibitem{bernstein} D. J. Bernstein, \emph{Cache-timing attacks on AES}, 2005.
\bibitem{russes} V. Arlazarov, E. Dinic, M. Kronrod, I. Faradzev, \emph{On economical construction of the transitive closure of an oriented graph}, Soviet Math.\ Doklady 11, 1970.
\bibitem{albrecht2011} M. R. Albrecht, G. V. Bard, C. Pernet, \emph{Efficient Dense Gaussian Elimination over the Finite Field with Two Elements}, arXiv:1111.6549, 2011.
\bibitem{mayoc} Dépôt officiel \code{PQCMayo/MAYO-C}, commit \code{74de637a7c4a7f9fafa888183bc05f5cf920cd2d}.
\end{thebibliography}

\end{document}
